\subsection{Perturbation of Riesz endomorphisms}

Lemma 1. --- Let E be a Banach space. Let p and q be continuous projectors on E such that \(\|q - p\| < 1\) . Then p induces an isomorphism of \(\operatorname{Im}(q)\) onto \(\operatorname{Im}(p)\) , and \(1_{E} - p\) induces an isomorphism of \(\operatorname{Ker}(q)\) onto \(\operatorname{Ker}(p)\) .

Consider the continuous linear maps \(u\colon x\mapsto p(x)\) from \(\operatorname{Im}(q)\) into \(\operatorname{Im}(p)\) and \(v\colon y\mapsto q(y)\) from \(\operatorname{Im}(p)\) into \(\operatorname{Im}(q)\) . For every \(x\in\operatorname{Im}(q)\) , we have \(x=q(x)\) , whence

\[
\begin{array}{c} (q - p) ^ {2} (x) = q ^ {2} (x) - q (p (x)) - p (q (x)) + p ^ {2} (x) \\ = x - q (p (x)) = x - v (u (x)). \end{array}
\]

It follows that we have \(\|1_{E}-v\circ u\|\leqslant\|q-p\|^{2}<1\) . By cor. 1 of I, p. 22, \(v\circ u\) is an automorphism of \(\operatorname{Im}(q)\) . One proves similarly that \(u\circ v\) is an automorphism of \(\operatorname{Im}(p)\) . This implies that u is an isomorphism from \(\operatorname{Im}(q)\) onto \(\operatorname{Im}(p)\) , whence the first assertion of the lemma. The second follows from it by replacing p and q by \(1_{E}-p\) and \(1_{E}-q\) respectively.

THEOREM 2. --- Let E be a Banach space. The set \(\mathcal{R}(\mathrm{E})\) of Riesz endomorphisms of E is open in \(\mathcal{L}(\mathrm{E})\) . The map that assigns to an element of \(\mathcal{R}(\mathrm{E})\) the dimension of its nilspace is upper semicontinuous on \(\mathcal{L}(\mathrm{E})\) .

The theorem follows from the following more precise statement:

PROPOSITION 5. --- Let E be a Banach space and \(u_{0}\) a Riesz endomorphism of E. Let N (resp. I) be the nilspace (resp. the conilspace) of \(u_{0}\) . Let d denote the dimension of N. There exists an open neighborhood U of \(u_{0}\) in \(\mathcal{L}(\mathrm{E})\) and an analytic mapping \(\pi\colon\mathrm{U}\to\mathcal{L}(\mathrm{E})\) such that

\begin{enumerate}
\def\labelenumi{(\roman{enumi})}
\item
  The endomorphism \(\pi(u_0)\) is the projector with image N and kernel I;
\item
  For every \(u \in U\) , the linear map \(\pi(u)\) is a projector of rank d that belongs to the bicommutant of u and, in particular, commutes with u. Moreover, u induces an automorphism of \(\operatorname{Ker}(\pi(u))\) ;
\item
  Every element of U is a Riesz endomorphism whose nilspace has dimension \(\leqslant d\) .
\end{enumerate}

If K = C, let \(\mathrm{Sp}(u_{0})\) denote the spectrum of \(u_{0}\) with respect to the algebra \(\mathcal{L}(\mathrm{E})\) . When \(K = R\) , let \(\mathrm{Sp}(u_{0})\) denote the complex spectrum \(\mathrm{Sp}_{\mathcal{L}(\mathrm{E}_{(\mathbf{C})})}((u_{0})_{(\mathbf{C})})\) of \(u_{0}\) (I, p. 85, no. 13). By Prop. 14 of III, p. 54, the point 0 is isolated in \(\{0\} \cup \mathrm{Sp}(u_{0})\) . Let r \textgreater{} 0 be a real number such that every element of \(\mathrm{Sp}(u_{0}) - \{0\}\) has modulus \textgreater{} r. Let V denote the open set in C consisting of complex numbers of absolute value ≠ r; let f be the holomorphic function on V defined by \(f(z) = 0\) if \(|z| > r\) and \(f(z) = 1\) if \(|z| < r\) . If K = C (resp. K = R), let U' denote the set of elements u of \(\mathcal{L}(\mathrm{E})\) whose spectrum (resp. whose complex spectrum) is contained in V. The set U' is an open neighborhood of \(u_{0}\) in \(\mathcal{L}(\mathrm{E})\) and the map \(u \mapsto f(u)\) from U' into \(\mathcal{L}(\mathrm{E})\) is holomorphic (I, p. 76, prop. 10 and I, p. 85, no. 13).

Let \(u \in \mathrm{U}'\) . The endomorphism \(f(u)\) is the spectral projector \(e_0(u)\) ; it commutes with \(u\) , and \(u\) induces, by passing to subspaces, an automorphism of \(\operatorname{Ker}(e_0(u))\) (cf. III, p. 53, no. 6).

The projector \(e_{0}(u_{0})\) has image N and kernel I (III, p. 54, prop. 14); its rank is d. By Lemma 1, the set U of elements \(u \in U'\) such that \(e_{0}(u)\) is of rank d is an open neighborhood of \(u_{0}\) . The set U and the map \(\pi: u \mapsto e_{0}(u)\) from U into \(\mathcal{L}(\mathrm{E})\) satisfy conditions (i) and (ii) of the proposition.

Let \(u \in U\) . Since u induces an automorphism of \(\operatorname{Ker}(e_{0}(u))\) , which is a closed vector subspace of finite codimension of E, the endomorphism u is a Riesz endomorphism of E (III, p. 48, prop. 8); the nilspace of u is contained in the image of \(\pi(u) = e_{0}(u)\) , and has dimension \(\leqslant d\) .

\subsection{Conorm of a continuous linear map}

Let E and F be normed spaces, and let \(u: E \rightarrow F\) be a continuous linear map, N the kernel of u and I its image. Let p denote the canonical surjection from E onto E/N and let i denote the canonical injection from I into F. Let \(\widetilde{u}\) be the bijective linear map from E/N onto I such that \(u = i \circ \widetilde{u} \circ p\) . The vector space E/N is equipped with the quotient norm, that is,

\[
\| y \| = \inf _ {x \in p^{-1}(y)} \| x \|\tag{2}
\]

for all \(y \in E/N\) . The map \(\widetilde{u}\) is continuous and \(\|u\| = \|\widetilde{u}\|\) , whence

\[
\| u\| = \sup_{\substack{y\in \mathrm{E / N}\\ y\neq 0}}\frac{\|\widetilde{u} (y)\|}{\|y\|},\tag{3}
\]

the supremum being taken in \(\overline{R}_{+}\) . The conorm of u is called the number

\[
((u)) = \inf_{\substack{y\in \mathrm{E} / \mathrm{N}\\ y\neq 0}}\frac{\| \widetilde{u}(y)\|}{\|y\|},\tag{4}
\]

the lower bound being taken in \(\overline{R}_{+}\) . We therefore have

\[
((u)) \| y \| \leqslant \| \widetilde {u} (y) \| \leqslant \| u \| \| y \|\tag{5}
\]

for every element \(y\) of E/N. Set \(v = i \circ \widetilde{u}\) . We have \(u = v \circ p\) and

\[
((u)) = \inf_{\substack{y\in \mathrm{E / N}\\ y\neq 0}}\frac{\|v(y)\|}{\|y\|} = \inf_{\substack{y\in \mathrm{E / N}\\ \|y\| = 1}}\| v(y)\| .\tag{6}
\]

When \(u\) is the zero map, the space E/N is reduced to 0 and one has \(((u)) = +\infty\) and \(\| u \| = 0\) . When \(u \neq 0\) , from (3) and (4) one deduces the inequalities

\[
0 \leqslant ((u)) \leqslant \| u \| <   + \infty .\tag{7}
\]

When u is injective, we have

\[
((u)) = \inf_{\substack{x\in \mathrm{E}\\ x\neq 0}}\frac{\|u(x)\|}{\|x\|},\tag{8}
\]

and, for all \(x \in \mathrm{E}\) , we have

\[
((u)) \| x \| \leqslant \| u (x) \| \leqslant \| u \| \| x \|.\tag{9}
\]

Denote by \(j\) the canonical injection of the normed space F into its completion \(\widehat{\mathbf{F}}\) . We have \(((u)) = ((j \circ u))\) .

Remark. --- By definition, \(((u)) > 0\) if and only if the bijective linear map \(\widetilde{u}\) is bicontinuous, that is, if and only if u is a strict morphism (TG, III, p. 16). We then have

\[
((u)) = \| \widetilde {u} ^ {- 1} \| ^ {- 1}.\tag{10}
\]

Lemma 2. --- Let c be a real number. If c \textless{} ((u)), then for every element \(z \in \operatorname{Im}(u)\) there exists an element x of E such that \(u(x) = z\) and \(c \|x\| \leqslant \|z\|\) . Conversely, if for every \(z \in \operatorname{Im}(u)\) there exists \(x \in E\) such that \(u(x) = z\) and \(c \|x\| \leqslant \|z\|\) , then \(c \leqslant ((u))\) .

This is a consequence of formulas (2) and (5) and of the definition of the conorm of \(u\) .

PROPOSITION 6. --- Let E and F be normed spaces and \(u \in \mathcal{L}(\mathrm{E};\mathrm{F})\) . Let B denote the set of elements of E with norm \(< 1\) . Set

\[
\mathrm{P} = u (\mathrm{E}) - u (\mathrm{B}), \quad \mathrm{Q} = u (\mathrm{E}) - (\overline {{u (\mathrm{B})}} \cap u (\mathrm{E})).
\]

The conorm of u is equal to the distance from 0 to P in F. If the normed space E is complete or if u is a strict morphism, the conorm of u is equal to the distance from 0 to Q in F.

Let N be the kernel of u; let \(p: E \to E/N\) be the canonical surjection, and let \(v: E/N \to F\) be the map induced by u on the quotient. The set \(p(B)\) is the set of elements of E/N of norm \textless{} 1. We have

\[
\mathrm{P} = u (\mathrm{E}) - u (\mathrm{B}) = v (\mathrm{E} / \mathrm{N}) - v (p (\mathrm{B})) = v ((\mathrm{E} / \mathrm{N}) - p (\mathrm{B}))
\]

since the map \(v\) is injective. In other words, P consists of the elements of F of the form \(v(y)\) with \(y \in \mathrm{E/N}\) and \(\|y\| \geqslant 1\) . Denote by \(d_{\mathrm{P}}\) the distance from 0 to P in F. We have

\[
d_{\mathrm{P}} = \inf_{\substack{y\in \mathrm{E} / \mathrm{N}\\ \| y\| \geqslant 1}}\| v(y)\| = \inf_{\substack{y\in \mathrm{E} / \mathrm{N}\\ \| y\| = 1}}\| v(y)\| = ((u))
\]

according to (6).

Suppose that \(u\) is a strict morphism. Let \(\varepsilon > 0\) . The set \(\varepsilon u(\mathrm{B})\) is a neighborhood of 0 in \(u(\mathrm{E})\) . The closure of \(u(\mathrm{B})\) in \(u(\mathrm{E})\) is equal to \(\overline{u(\mathrm{B})} \cap u(\mathrm{E})\) . It is contained in the set \(u(\mathrm{B}) + \varepsilon u(\mathrm{B})\) which is equal to \((1 + \varepsilon)u(\mathrm{B})\) since \(u(\mathrm{B})\) is convex. Consequently, we have \((1 + \varepsilon)\mathrm{P} \subset \mathrm{Q} \subset \mathrm{P}\) and the distance \(d_{\mathrm{Q}}\) from 0 to Q in F satisfies the inequalities \(d_{\mathrm{P}} \leqslant d_{\mathrm{Q}} \leqslant (1 + \varepsilon)d_{\mathrm{P}}\) . Since this holds for every \(\varepsilon > 0\) , we have \(d_{\mathrm{Q}} = d_{\mathrm{P}} = ((u))\) .

Suppose that \(u\) is not a strict morphism, but that the normed space E is complete. We then have \(((u)) = 0\) . The closure of \(u(\mathrm{B})\) in \(u(\mathrm{E})\) , which is equal to \(\overline{u(\mathrm{B})} \cap u(\mathrm{E})\) , is not a neighborhood of 0 in \(u(\mathrm{E})\) (EVT, I, p. 17, th. 1). There therefore exist points of Q arbitrarily close to 0, whence \(d_{\mathrm{Q}} = 0 = ((u))\) .

PROPOSITION 7. --- Let E be a Banach space and F a normed space. The mapping \(u \mapsto ((u))\) of \(\mathcal{L}(\mathrm{E};\mathrm{F})\) in \(\overline{\mathbf{R}}\) is upper semi-continuous.

Let \(u \in \mathcal{L}(\mathrm{E};\mathrm{F})\) . It suffices to prove that for every real number \(c > ((u))\) , the set of elements \(v \in \mathcal{L}(\mathrm{E};\mathrm{F})\) such that \(((v)) < c\) is a neighborhood of \(u\) . Let B denote the set of elements of E with norm \(< 1\) . By prop. 6, there exists \(y \in \mathrm{E}\) such that \(u(y) \notin \overline{u(\mathrm{B})}\) and \(\| u(y)\| < c\) . The distance \(d\) of \(u(y)\) to the closed set \(\overline{u(\mathrm{B})}\) is strictly positive. The set V of elements \(v\) of \(\mathcal{L}(\mathrm{E};\mathrm{F})\) satisfying the relations \(\| v(y)\| < c\) and \(\| u - v\|(1 + \| y\|) < d\) is a neighborhood of \(u\) in \(\mathcal{L}(\mathrm{E};\mathrm{F})\) . Let \(v \in \mathrm{V}\) . For every \(x \in \mathrm{B}\) , we have

\[
\begin{array}{c} d \leqslant \| u (y) - u (x) \| \leqslant \| v (y) - v (x) \| + \| u - v \| (\| y \| + \| x \|) \\ <   \| v (y) - v (x) \| + d. \end{array}
\]

Consequently \(v(y)\) does not belong to \(v(\mathrm{B})\) , and one has \(((v)) \leqslant \| v(y) \| < c\) by prop. 6.

PROPOSITION 8. --- Let E be a Banach space, F a normed space. For every \(u \in \mathcal{L}(\mathrm{E};\mathrm{F})\) , we have \(((u)) = ((^t u))\) .

Denote by \(j\) the canonical injection of the normed space F into its completion \(\widehat{\mathrm{F}}\) . We have \(((u)) = ((j \circ u))\) . Since the linear map \(^t j: (\widehat{\mathrm{F}})' \to \mathrm{F}'\) is bijective and isometric, one also has \(((^t u)) = ((^t (j \circ u)))\) . It therefore suffices to prove the proposition when the normed space F is complete, which we assume in the rest of the proof.

If \(u\) is zero, one has \(((u)) = ((^t u)) = +\infty\) . If \(u\) is not a strict morphism, then \(^t u\) is not one either (EVT, IV, p. 29, cor. 3), and one has \(((u)) = ((^t u)) = 0\) .

Suppose now that u is a nonzero strict morphism. Let N denote the kernel of u and I its image, and consider the canonical decomposition of u:

\[
\mathrm{E} \xrightarrow {p} \mathrm{E/N} \xrightarrow {v} \mathrm{I} \xrightarrow {i} \mathrm{F}.
\]

The kernel of \(^{t}u\) is the orthogonal I° of I in F' (EVT, IV, p. 27, prop. 2), and \(^{t}i\) induces by passage to the quotient an isometry \(\iota\) from F'/I° onto I' (EVT, IV, p. 9, prop. 10). Moreover, \(^{t}p\) defines an isometry from (E/N)' onto the orthogonal N° of N in E' (EVT, IV, p. 8, prop. 9). The canonical decomposition of \(^{t}u\) is therefore

\[
\mathrm{F} ^ {\prime} \longrightarrow \mathrm{F} ^ {\prime} / \mathrm{I} ^ {\circ} \stackrel {v ^ {\prime}} {\longrightarrow} \mathrm{N} ^ {\circ} \longrightarrow \mathrm{E} ^ {\prime},
\]

where \(v^{\prime} = \pi \circ{}^{t}v\circ \iota\). We then have

\[
\| (v ^ {\prime}) ^ {- 1} \| = \| (^ {t} v) ^ {- 1} \| = \| ^ {t} (v ^ {- 1}) \| = \| v ^ {- 1} \|
\]

(EVT, IV, p. 7, prop. 8), whence \(((u)) = ((^t u))\) according to formula (10).

\subsection{Finite-dimensional vector subspaces of a normed space}

The following statement will be proved in TA, forthcoming, as a corollary of the Borsuk–Ulam theorem.

THEOREM 3. — Let n be a positive integer, V a real normed vector space of dimension \(n + 1\) and W a vector subspace of V of dimension n. Let S be the unit sphere of V. There does not exist a continuous map \(f: S \to W\) such that \(\|f(x) - x\| < 1\) for all \(x \in S\) .

THEOREM 4 (Krein, Krasnoselskii, Milman). — Let E be a normed space, and let F and G be vector subspaces of E. Suppose that the dimension of G is finite and strictly less than that of F. There exists an element of F of norm 1 whose distance to G is equal to 1.

It suffices to treat the case where the field K is equal to R. Let n be the dimension of G. By replacing F by a vector subspace of F of dimension \(n+1\) containing G, one reduces to the case where \(\dim(F)=n+1\) . Argue by contradiction by assuming that the conclusion of the theorem is not satisfied. Let S be the unit sphere of F. For every \(x\in S\) , the distance from x to G is then strictly less than \(\|x\|=1\) and one can choose an element \(y(x)\) of G for which one has \(\|x-y(x)\|<1\) . Denote by \(V_{x}\) the set of elements z of S such that \(\|z-y(x)\|<1\) ; this is an open neighborhood of x in S. There exists a locally finite continuous partition of unity \((\varphi_{x})_{x\in\mathrm{S}}\) on S subordinate to the covering \((\mathrm{V}_{x})_{x\in\mathrm{S}}\) of S (TG, IX, p. 46, prop. 3 and p. 51, prop. 6). Let \(f\colon S\to G\) be the continuous map given by

\[
f (z) = \sum_ {x \in \mathrm{S}} \varphi_ {x} (z) y (x).
\]

For all \(z \in S\), we have \(\varphi_x(z) \geqslant 0\) for all \(x \in S\), and there exists \(x \in S\) such that \(\varphi_x(z) > 0\) since \(\sum_{x \in S} \varphi_x(z) = 1\), hence

\[
\begin{array}{c} \| z - f (z) \| = \left\| \sum_ {x \in \mathrm{S}} \varphi_ {x} (z) (z - y (x)) \right\| \leqslant \sum_ {x \in \mathrm{S}} \varphi_ {x} (z) \| z - y (x) \| \\ <   \sum_ {x \in \mathrm{S}} \varphi_ {x} (z) = 1. \end{array}
\]

This property of f contradicts Theorem 3.

PROPOSITION 9. — Let E and F be normed spaces, \(n \in N\) and let u be a continuous linear map from E into F whose kernel has dimension n. The conorm of u is equal to the distance d from u to the set of maps \(v \in \mathcal{L}(\mathrm{E};\mathrm{F})\) whose kernel has dimension at least \(n + 1\) .

When \(u = 0\) , we have \(((u)) = +\infty\) and \(\dim(\mathrm{E}) = n\) , whence \(d = +\infty\) . Suppose henceforth that \(u\) is nonzero, and therefore \(((u)) < +\infty\) . Let \(v\) be an element of \(\mathcal{L}(\mathrm{E};\mathrm{F})\) such that \(\| u - v \| < ((u))\) and let us prove that its kernel has dimension \(\leqslant n\) . Let \(x\) be an element of norm 1 of \(\operatorname{Ker}(v)\) and \(y\) be its image in \(\operatorname{E}/\operatorname{Ker}(u)\) . We have (formula (5) of III, p. 61)

\[
((u)) \| y \| \leqslant \| u (x) \| = \| (u - v) (x) \| \leqslant \| u - v \| <   ((u))
\]

whence \(\|y\| < 1\) . Now \(\|y\|\) is the distance from x to \(\operatorname{Ker}(u)\) . Theorem 4 then implies that we have \(\dim \operatorname{Ker}(v) \leqslant \dim \operatorname{Ker}(u) = n\) . This proves the inequality \(((u)) \leqslant d\) . The reverse inequality \(d \leqslant ((u))\) follows from the more precise lemma below.

Lemma 3. — Let c be a real number such that \(c > ((u))\) . There exists a continuous linear map \(h: E \to F\) , of rank 1 and of norm \textless{} c such that the kernel of \(u + h\) contains that of u and is of dimension \(n + 1\) .

Denote by \(p\) the canonical mapping of E onto \(\mathrm{E} / \mathrm{Ker}(u)\) . There exists \(a \in \mathrm{E}\) such that \(\| p(a) \| = 1\) and \(\| u(a) \| < c\) (formula (6) of III, p. 61). By the Hahn--Banach theorem (EVT, II, p. 24, cor. 2), there exists a continuous linear form \(f\) on the normed space \(\mathrm{E} / \mathrm{Ker}(u)\) such that \(f(p(a)) = 1\) and \(\| f \| = 1\) . The linear map \(h: x \mapsto -(f \circ p)(x)u(a)\) from E into F is continuous, of rank 1, and we have \(\| h \| \leqslant \| f \| \| p \| \| u(a) \| < c\) . The kernel of the map \(u + h\) contains that of \(u\) and \(a\) ; since \(a \notin \mathrm{Ker}(u)\) , its dimension is at least \(n + 1\) . On the other hand, the linear map \(u\) induces, by passing to subspaces, a linear map from \(\mathrm{Ker}(u + h)\) into \(\mathrm{Im}(h)\) with kernel \(\mathrm{Ker}(u) \cap \mathrm{Ker}(u + h)\) , so that \(\dim(\mathrm{Ker}(u + h)) \leqslant \dim(\mathrm{Ker}(u)) + 1\) . The conclusion follows.

COROLLARY 1. --- Let E and F be normed spaces, and let u, v be nonzero continuous linear maps from E into F whose kernels have the same finite dimension. Then we have

\[
\left| \left((u)\right) - ((v)) \right| \leqslant \| u - v \|.
\]

Denote by \(n\) the common dimension of the kernels of \(u\) and \(v\). Let A be the set of continuous linear maps from E to F whose kernel has dimension \(\geqslant n + 1\) . Since \(u \neq 0\) , we have \(\dim(\mathrm{E}) > n\) and the set A contains 0. By Proposition 9, \(((u))\) and \(((v))\) are respectively the distances from \(u\) and \(v\) to the set A in \(\mathcal{L}(\mathrm{E};\mathrm{F})\). It then suffices to apply the formula \(|d(u,\mathrm{A}) - d(v,\mathrm{A})| \leqslant \|u - v\|\) (cf. TG, IX, p. 13).

COROLLARY 2. --- Let E and F be Banach spaces, and let u and v be nonzero continuous linear maps from E into F whose images have the same finite codimension in F. Then one has

\[
| ((u)) - ((v)) | \leqslant \| u - v \|.
\]

The morphism u is strict (III, p. 52, lemma 6). The kernel of the continuous linear map \({}^{t}u\) is the orthogonal \((\mathrm{Im}(u))^{\circ}\) of \(\mathrm{Im}(u)\) (EVT, IV, p. 27, prop. 2), hence \(\dim(\mathrm{Ker}({}^{t}u)) = \mathrm{codim}_{\mathrm{F}}(\mathrm{Im}(u))\) , and similarly \(\dim(\mathrm{Ker}({}^{t}v)) = \mathrm{codim}_{\mathrm{F}}(\mathrm{Im}(v))\) . The kernels of \({}^{t}u\) and \({}^{t}v\) therefore have the same finite dimension. Since

\[
\| ^ {t} u - ^ {t} v \| = \| u - v \|, \quad ((^ {t} u)) = ((u)), \quad ((^ {t} v)) = ((v))
\]

(EVT, IV, p. 7, prop. 8 and prop. 8 of III, p. 63), the assertion follows from corollary 1, applied to \(^{t}u\) and \(^{t}v\) .

\subsection{Perturbations of continuous injective or surjective linear maps}

In this section, we adopt the following conventions: if E is a finite-dimensional vector space, \(\dim(\mathrm{E})\) denotes its dimension; if E is an infinite-dimensional vector space, we set \(\dim(\mathrm{E}) = +\infty \in \overline{\mathbf{R}}\) . If u is a linear map whose kernel or cokernel is finite-dimensional, we set \(\operatorname{ind}(u) = \dim \operatorname{Coker}(u) - \dim \operatorname{Ker}(u)\) , the computation being carried out in \(\overline{R}\) .

Let E and F be normed spaces. We denote by \(\mathcal{M}(\mathrm{E};\mathrm{F})\) the set of strict injective morphisms from E into F, and by \(\mathcal{Q}\mathcal{M}(\mathrm{E};\mathrm{F})\) the set of strict morphisms from E into F whose kernel is finite-dimensional.

PROPOSITION 10. --- Let E and F be normed spaces. The set \(\mathcal{M}(\mathrm{E};\mathrm{F})\) is open in \(\mathcal{L}(\mathrm{E};\mathrm{F})\) . It is the interior of the set of maps in \(\mathcal{L}(\mathrm{E};\mathrm{F})\) that are injective.

Let A be the set of continuous injective linear maps from E into F. For a map \(u \in A\) to be a strict morphism, it is necessary and sufficient that its conorm \(((u))\) be strictly positive (III, p. 62, remark). Now \(((u))\) is the distance from \(u\) to the complement of A in \(\mathcal{L}(\mathrm{E};\mathrm{F})\) (III, p. 65, prop. 9). The proposition follows.

PROPOSITION 11. --- Let E and F be Banach spaces. The set \(\mathcal{Q}\mathcal{M}(\mathrm{E};\mathrm{F})\) is open in \(\mathcal{L}(\mathrm{E};\mathrm{F})\) . It is the interior of the set of maps in \(\mathcal{L}(\mathrm{E};\mathrm{F})\) whose kernel is finite-dimensional.

Let A be the set of continuous linear maps from E into F whose kernel is finite-dimensional.

Let u be an element of \(\mathcal{Q}\mathcal{M}(\mathrm{E};\mathrm{F})\) . We then have \(((u)) > 0\) (III, p. 62, remark). Every element \(v \in \mathcal{L}(\mathrm{E};\mathrm{F})\) whose distance from u is \(<((u))\) then belongs to A (III, p. 65, prop. 9), so that u is an interior point of A.

Conversely, let \(u\) be an element of A that is not a strict morphism. We have \(((u)) = 0\) (III, p. 62, remark). Let \(v\) be an element of \(\mathcal{L}(\mathrm{E};\mathrm{F})\) which differs from \(u\) by a linear map of finite rank; by Corollary 1 of III, p. 40, the morphism \(v\) cannot be strict with closed image; since a strict morphism from E into F has a closed image in F (EVT, IV, p. 28, th. 1), this means that v is not a strict morphism. We therefore have \(((v)) = 0\) . Let \(\varepsilon\) be a real number \textgreater{} 0. By the preceding observations, Lemma 3 of III, p. 66 allows us to construct by induction a sequence \((u_{m})_{m \in \mathbf{N}}\) of elements of \(\mathcal{L}(\mathrm{E}; \mathrm{F})\) satisfying the following conditions:

\begin{enumerate}
\def\labelenumi{(\roman{enumi})}
\item
  One has \(u_{0}=u\) ;
\item
  For every \(m \geqslant 0\) , \(u_{m+1} - u_m\) is a continuous linear map of rank 1 and norm \(\leqslant 2^{-m-1}\varepsilon\) ;
\item
  For every \(m \geqslant 0\) , the kernel of \(u_m\) has dimension \(n + m\) and is contained in that of \(u_{m+1}\) .
\end{enumerate}

The sequence \((u_{m})\) is a Cauchy sequence in the Banach space \(\mathcal{L}(\mathrm{E};\mathrm{F})\) . Let \(u'\) be its limit. The kernel of \(u'\) contains that of \(u_{m}\) for all \(m \geqslant 0\) ; it is of infinite dimension. Since

\[
\left\| u ^ {\prime} - u \right\| \leqslant \sum_ {m = 0} ^ {\infty} \left\| u _ {m + 1} - u _ {m} \right\| \leqslant \varepsilon \sum_ {m = 0} ^ {\infty} 2 ^ {- m - 1} = \varepsilon ,
\]

the distance from u to the complement of A is less than \(\varepsilon\) . Since this holds for every \(\varepsilon > 0\) , one concludes that u is not an interior point of A.

PROPOSITION 12. --- Let E and F be Banach spaces and u an element of \(\mathcal{Q}\mathcal{M}(\mathrm{E};\mathrm{F})\) . Every \(v \in \mathcal{L}(\mathrm{E};\mathrm{F})\) such that \(\|v - u\| < ((u))\) belongs to \(\mathcal{Q}\mathcal{M}(\mathrm{E};\mathrm{F})\) and satisfies the relations

\[
\begin{array}{c} \dim \operatorname{Ker} (v) \leqslant \dim \operatorname{Ker} (u), \\ \dim \operatorname{Coker} (v) \leqslant \dim \operatorname{Coker} (u), \\ \operatorname{ind} (v) = \operatorname{ind} (u). \end{array}
\]

Since \(u\) is strict, we have \(((u)) > 0\) . Let B denote the open ball with center \(v\) and radius \(((u))\) in \(\mathcal{L}(\mathrm{E};\mathrm{F})\) . For every \(v \in \mathrm{B}\) , we have \(\dim \operatorname{Ker}(v) \leqslant \dim \operatorname{Ker}(u)\) (III, p. 65, prop. 9) and \(v \in \mathcal{Q}\mathcal{M}(\mathrm{E};\mathrm{F})\) (prop. 11).

For \(r \in \mathbf{Z} \cup \{+\infty\}\) , denote by \(B_r\) the set of elements \(v \in B\) such that \(\text{ind}(v) = r\) . If \(r \in \mathbf{Z}\) , the set \(B_r\) is the set of Fredholm maps from E to F of index \(r\) which belong to B (III, p. 52, prop. 11); it is open in B by th. 1 of III, p. 58.

Let us prove that the sets \(B_{r}\) are closed in B. Let \(v \in B\) be a point adherent to \(B_{r}\) . Since the sets \(B_{s}\) , for \(s \in Z\) , are open in B and pairwise disjoint, we have \(v \in B_{r}\) or \(v \in B_{+\infty}\) .

Suppose \(v \in B_{+\infty}\) . Let n denote the dimension of \(\operatorname{Ker}(u)\) . Choose a vector subspace T of F of dimension \(r + 2n + 1\) whose intersection with \(\operatorname{Im}(v)\) is reduced to 0 and a topological supplement S of \(\operatorname{Ker}(v)\) in E (III, p. 55, prop. 1). Consider the continuous linear map \(f: (s, t) \mapsto v(s) + t\) from S × T to F. It is injective. The linear map v is a strict morphism since v belongs to B ⊂ \(\mathcal{Q}\mathcal{M}(\mathrm{E};\mathrm{F})\) . The image of v is therefore closed (EVT, IV, p. 28, th. 1). By prop. 1 of III, p. 39, the restriction of v to S is a strict morphism with closed image from S to F, and f is a strict morphism with closed image from S × T to F.

By prop. 10, there exists a neighborhood U of \(v\) in B such that, for every \(w \in \mathrm{U}\) , the linear map \((s,t) \mapsto w(s) + t\) from \(\mathrm{S} \times \mathrm{T}\) into F is injective. Let \(w\) be an element of U. We then have \(\operatorname{Ker}(w) \cap \mathrm{S} = \{0\}\) , hence \(w\) induces by restriction and passage to the quotient an injective linear map from \(\operatorname{Ker}(w)\) into E/S, which implies that \(\operatorname{Ker}(w)\) has dimension at most \(n\) . We also have \(w(\mathrm{S}) \cap \mathrm{T} = \{0\}\) , hence

\[
\operatorname{codim} _ {\mathrm{F}} (\operatorname{Im} (w)) \geqslant \dim (\mathrm{T}) - \operatorname{codim} _ {\mathrm{E}} (\mathrm{S}) = r + n + 1,
\]

this entails \(\operatorname{ind}(w) \geqslant r + 1\) and contradicts the hypothesis that v is adherent to B \(_{r}\) .

If \(r\) is an element of \(\mathbf{Z}\) such that \(B_r\) is nonempty, we have \(B_r = B\) since \(B_r\) is open and closed in \(B\) and \(B\) is connected. If \(B_r\) is empty for every \(r \in \mathbf{Z}\) , we have \(\text{ind}(v) = +\infty\) for all \(v \in B\) . The map \(v \mapsto \text{ind}(v)\) is therefore constant on \(B\) . Finally, for every \(v \in B\) , we have

\[
\begin{array}{c} \dim \operatorname{Coker} (v) = \operatorname{ind} (v) + \dim \operatorname{Ker} (v) \\ \leqslant \operatorname{ind} (u) + \dim \operatorname{Ker} (u) = \dim \operatorname{Coker} (u). \end{array}
\]

This concludes the proof.

COROLLARY 1. --- The functions defined by \(u \mapsto \dim \operatorname{Ker}(u)\) and by \(u \mapsto \dim \operatorname{Coker}(u)\) on \(\mathcal{Q}\mathcal{M}(\mathrm{E};\mathrm{F})\) are upper semicontinuous. The function \(u \mapsto \operatorname{ind}(u)\) is locally constant on \(\mathcal{Q}\mathcal{M}(\mathrm{E};\mathrm{F})\) .

COROLLARY 2. --- The function \(u \mapsto \dim \operatorname{Coker}(u)\) is locally constant on the set \(\mathcal{M}(\mathrm{E};\mathrm{F})\) of strict injective morphisms from E into F.

Indeed, one has \(\dim \operatorname{Coker}(u) = \operatorname{ind}(u)\) for \(u \in \mathcal{M}(\mathrm{E};\mathrm{F})\) .

Lemma 4. --- Let E and F be Banach spaces. For an element u of \(\mathcal{L}(\mathrm{E};\mathrm{F})\) to be a strict morphism from E into F, it is necessary and sufficient that \(^{t}u\) be a strict morphism from F' into E'. In this case, one has \(\dim \operatorname{Coker}(^{t}u)=\dim\operatorname{Ker}(u)\) and \(\dim\operatorname{Ker}(^{t}u)=\dim\operatorname{Coker}(u)\) .

For u to be a strict morphism, it is necessary and sufficient that \(^{t}u\) be one (EVT, IV, p. 29, cor. 3); in this case, the image of u is closed (EVT, IV, p. 28, th. 1) and the vector space \(\mathrm{Ker}(^{t}u)\) (resp. \(\mathrm{Coker}(^{t}u)\) ) is canonically isomorphic to the dual of the normed space \(\mathrm{Coker}(u)\) (resp. \(\mathrm{Ker}(u)\) ) according to EVT, IV, p. 27, prop. 2. The lemma follows.

Let E and F be Banach spaces. We denote by \(\mathcal{E}(\mathrm{E};\mathrm{F})\) the set of continuous surjective linear maps from E into F and by \(\mathcal{Q}\mathcal{E}(\mathrm{E};\mathrm{F})\) the set of continuous linear maps from E into F whose image has finite codimension. Every element of \(\mathcal{Q}\mathcal{E}(\mathrm{E};\mathrm{F})\) is a strict morphism with closed image (III, p. 52, lemma 6). It follows from lemma 4 that \(\mathcal{E}(\mathrm{E};\mathrm{F})\) and \(\mathcal{Q}\mathcal{E}(\mathrm{E};\mathrm{F})\) are respectively the inverse images of \(\mathcal{M}(\mathrm{F}^{\prime};\mathrm{E}^{\prime})\) and \(\mathcal{Q}\mathcal{M}(\mathrm{F}^{\prime};\mathrm{E}^{\prime})\) under the continuous map \(u\mapsto{}^{t}u\) from \(\mathcal{L}(\mathrm{E};\mathrm{F})\) into \(\mathcal{L}(\mathrm{F}^{\prime};\mathrm{E}^{\prime})\) .

PROPOSITION 13. --- Let E and F be Banach spaces. The sets \(\mathcal{E}(\mathrm{E};\mathrm{F})\) and \(\mathcal{Q}\mathcal{E}(\mathrm{E};\mathrm{F})\) are open in \(\mathcal{L}(\mathrm{E};\mathrm{F})\) . More precisely, if \(u\) is an element of \(\mathcal{Q}\mathcal{E}(\mathrm{E};\mathrm{F})\) and \(v\) is an element of \(\mathcal{L}(\mathrm{E};\mathrm{F})\) such that \(\| v - u\| < ((u))\) , we have \(v \in \mathcal{Q}\mathcal{E}(\mathrm{E};\mathrm{F})\) and

\(\dim \operatorname{Ker}(v) \leqslant \dim \operatorname{Ker}(u), \quad \dim \operatorname{Coker}(v) \leqslant \dim \operatorname{Coker}(u),\)

\[
\operatorname{ind} (v) = \operatorname{ind} (u).
\]

We saw above that \(^{t}u \in \mathcal{Q}\mathcal{M}(F';E')\) . Moreover, for every element \(v\) of \(\mathcal{L}(E;F)\) , we have \(\|^{t}v - ^{t}u\| = \|v - u\|\) and \(((^{t}u)) = ((u))\) (EVT, IV, p. 7, prop. 8 and III, p. 63, prop. 8). By prop. 12, it follows from these relations that if \(\|v - u\| < ((u))\) , then \(^{t}v\) belongs to \(\mathcal{Q}\mathcal{M}(F';E')\) and one has the inequalities

\[
\dim \operatorname{Ker} (^ {t} v) \leqslant \dim \operatorname{Ker} (^ {t} u), \quad \dim \operatorname{Coker} (^ {t} v) \leqslant \dim \operatorname{Coker} (^ {t} u)
\]

as well as the equality \(\mathrm{ind}(^{t}v)=\mathrm{ind}(^{t}u)\) . The proposition then follows from lemma 4.

COROLLARY 1. --- The functions defined by \(u \mapsto \dim \operatorname{Ker}(u)\) and by \(u \mapsto \dim \operatorname{Coker}(u)\) on \(\mathcal{Q}\mathcal{E}(\mathrm{E};\mathrm{F})\) are upper semicontinuous. The function \(u \mapsto \operatorname{ind}(u)\) is locally constant on \(\mathcal{Q}\mathcal{E}(\mathrm{E};\mathrm{F})\) .

COROLLARY 2. --- The function defined by \(u \mapsto \dim \operatorname{Ker}(u)\) is locally constant on \(\mathcal{E}(\mathrm{E};\mathrm{F})\) .

Indeed, one has \(\dim \operatorname{Ker}(u) = -\operatorname{ind}(u)\) for all \(u \in \mathcal{E}(\mathrm{E};\mathrm{F})\) .

\section{PERTURBATION BY A COMPACT LINEAR MAP}

\subsection{Strict morphisms and properness}

Lemma 1. --- Let E and F be topological vector spaces, let u be a continuous linear map from E into F, and let U be a neighborhood of 0 in E. Suppose that u induces a homeomorphism from U onto a closed subset of F. Then the image I of u is closed and u induces a homeomorphism from E onto I.

The set \(u(\mathrm{U})\) is a neighborhood of 0 in I; it is closed in F, hence the subgroup I of F is locally closed at 0, and consequently it is closed (TG, III, p. 7, prop. 4).

Since \(\operatorname{Ker}(u)\) meets U only at 0, one has \(\operatorname{Ker}(u) = \{0\}\) and the map \(u\) is injective. Let V be a closed neighborhood of 0 in E contained in U. Since \(u(\mathring{\mathrm{V}})\) is a neighborhood of 0 in \(u(\mathrm{V})\) , there exists a balanced neighborhood W of 0 in F such that \(u(\mathrm{V}) \cap \mathrm{W} \subset u(\mathring{\mathrm{V}})\) . The set \(u^{-1}(\mathrm{W})\) is balanced, hence connected. It contains 0 and is contained in \(\mathring{\mathrm{V}} \cup (\mathrm{E} - \mathrm{V})\) . Since \(\mathring{\mathrm{V}}\) and E-V are disjoint open subsets of E, the set \(u^{-1}(\mathrm{W})\) is contained in \(\mathring{\mathrm{V}}\) , whence \(\mathrm{W} \cap \mathrm{I} \subset u(\mathring{\mathrm{V}})\) . Consequently \(u(\mathring{\mathrm{V}})\) is a neighborhood of 0 in I. This implies that \(u\) induces a homeomorphism from E onto I.

PROPOSITION 1. --- Let E be a separated locally convex space, F a locally convex space and u a continuous linear map from E into F. The following conditions are equivalent:

\begin{enumerate}
\def\labelenumi{(\roman{enumi})}
\item
  The map u is a strict morphism, its kernel has finite dimension and its image is closed in F;
\item
  There exists a closed neighborhood V of 0 in E such that the restriction of u to V is proper (TG, I, p. 72).
\item
  \(\Longrightarrow\) (ii): Suppose that u satisfies condition (i). Since the kernel of u has finite dimension, it has a topological supplement \(E_{1}\) in E (III, p. 55, prop. 1), and there exists a compact neighborhood C of 0 in \(\operatorname{Ker}(u)\) . Identify E with \(E_{1} \times \operatorname{Ker}(u)\) ; the set \(V = E_{1} \times C\) is then a closed neighborhood of 0 in E. The restriction of u to V is the composition of the projection of \(E_{1} \times C\) onto \(E_{1}\) which is proper (TG, I, p. 77, cor. 5) and of the restriction \(u_{1}\) of u to \(E_{1}\) . Now \(u_{1}\) is a homeomorphism from \(E_{1}\) onto a closed subspace of F, hence is proper (TG, I, p. 72, prop. 2). The composition of two proper maps is proper (TG, I, p. 73, prop. 5, a)), so the restriction of u to V is proper.
\item
  \(\Longrightarrow\) (i): Let V be a closed neighborhood of 0 in E such that the restriction v of u to V is proper. The set \(\mathrm{V} \cap \mathrm{Ker}(u) = \overline{v}(\{0\})\) is then compact (TG, I, p. 75, th. 1); consequently, the vector space \(\mathrm{Ker}(u)\) is locally compact, hence finite-dimensional (EVT, I, p. 15, th. 3). Let \(E_{1}\) be a topological complement of \(\mathrm{Ker}(u)\) in E (prop. 1 of III, p. 55); set \(V_{1} = E_{1} \cap V\) . The set \(V_{1}\) is closed in V. The map \(u|V_{1}\) is proper (TG, I, p. 74, cor. 1) and injective, hence is a homeomorphism from \(V_{1}\) onto a closed subset of F (TG, I, p. 72, prop. 2). By Lemma 1, the restriction of u to \(E_{1}\) is a homeomorphism from \(E_{1}\) onto a closed subspace of F, hence u is a strict morphism with closed image.
\end{enumerate}

\subsection{Perturbation of injective or surjective linear maps}

THEOREM 1. --- Let E and F be separated locally convex spaces and let u be a strict morphism from E into F, whose kernel is finite-dimensional and whose image is closed. Let h be a compact linear map from E into F. The linear map \(u + h\) is a strict morphism, its kernel is finite-dimensional, and its image is closed.

By prop. 1 of III, p. 71, there exists a closed neighborhood V of 0 in E such that the restriction of \(u\) to V is proper. Since \(h\) is compact, there exists a closed neighborhood W of 0 contained in V such that the set \(h(\mathrm{W})\) is relatively compact. Set \(\mathrm{C} = \overline{h(\mathrm{W})}\) . The restriction of \(u\) to W is proper (TG, I, p. 74, cor. 1). The map \(\alpha: x \mapsto (u(x), h(x))\) from W into \(\mathrm{F} \times \mathrm{C}\) is proper because the map \(\mathrm{pr}_1 \circ \alpha = u|\mathrm{W}\) from W into F is proper (TG, I, p. 73, prop. 5, \(d\) ). The map \(\beta: (y, z) \mapsto (y + z, z)\) from \(\mathrm{F} \times \mathrm{C}\) into \(\mathrm{F} \times \mathrm{C}\) is a homeomorphism, and the map \(\mathrm{pr}_1\) from \(\mathrm{F} \times \mathrm{C}\) into F is proper (TG, I, p. 77, cor. 5). The composite map \(\mathrm{pr}_1 \circ \beta \circ \alpha\) from W into F is therefore proper (TG, I, p. 73, prop. 5, \(a\) ). But this map is none other than the restriction of \(u + h\) to W. By prop. 1 of III, p. 71, \(u + h\) is a strict morphism, its kernel is finite-dimensional, and its image is closed.

Lemma 2. --- Let E and F be Fréchet spaces and \(u \in \mathcal{L}(\mathrm{E};\mathrm{F})\) . The following conditions are equivalent:

\begin{enumerate}
\def\labelenumi{(\roman{enumi})}
\item
  The cokernel of u is finite-dimensional;
\item
  \(^{t}u \colon \mathrm{F}_c' \to \mathrm{E}_c'\) is a strict morphism with closed image whose kernel is finite-dimensional.
\item
  \(\Longrightarrow\) (ii): Suppose that the cokernel of u is finite-dimensional. The map u is then a strict morphism (Lemma 6 of III, p. 52). By Prop. 2 of EVT, IV, p. 27, the image of \(^{t}u\) is closed in E' equipped with the weak topology, and a fortiori in \(E_{c}'\) . The kernel of \(^{t}u\) is the orthogonal complement of the image of u (loc. cit.); it is therefore finite-dimensional. Finally, \(^{t}u\) is a strict morphism from \(F_{c}'\) into \(E_{c}'\) by Theorem 1 of EVT, IV, p. 28.
\item
  \(\Longrightarrow\) (i): Suppose that \(^t u\colon F_c' \to E_c'\) is a strict morphism. By EVT, IV, p. 28, th. 1, the image of \(u\) is closed. The kernel of \(^t u\) is the orthogonal complement of \(\operatorname{Im}(u)\) (by EVT, IV, p. 27, prop. 2); if this kernel is finite-dimensional, the image of \(u\) is of finite codimension in F.
\end{enumerate}

THEOREM 2. --- Let E and F be Fréchet spaces, u: E → F a continuous linear map whose cokernel is finite-dimensional, and h: E → F a compact linear map. The linear map u + h is a strict morphism, its cokernel is finite-dimensional, and its image is closed.

By lemma 6 of III, p. 52, it suffices to show that the cokernel of \(u + h\) is of finite dimension. Now, by prop. 9 of III, p. 9, the map \(^{t}h\) from \(F_{c}^{\prime}\) into \(E_{c}^{\prime}\) is compact. Theorem 2 follows from Theorem 1 and Lemma 2.

\subsection{Perturbation of Fredholm mappings}

THEOREM 3. --- Let E be a locally convex space, F a separated locally convex space, u a Fredholm map from E into F, and h a compact linear map from E into F. Then \(u + h\) is a Fredholm operator, and one has \(\operatorname{ind}(u + h) = \operatorname{ind}(u)\) .

We will prove the theorem in three steps.

\begin{enumerate}
\def\labelenumi{\Alph{enumi})}
\item
  Assume that E and F are Banach spaces. Prop. 2 of III, p. 42 shows that the linear map u is a strict morphism with closed image whose kernel and cokernel are finite-dimensional. Since E and F are Banach spaces, Ths. 1 and 2 of III, p. 73 imply that, for every \(t \in [0,1]\) , the linear map \(u_t = u + th\) possesses these same properties, and is therefore a Fredholm operator (III, p. 42, prop. 2). The map \(t \mapsto u_t\) from \([0,1]\) to the set \(\mathcal{F}(\mathrm{E};\mathrm{F})\) of Fredholm maps from E into F is continuous. By th. 1 of III, p. 58, the map \(t \mapsto \operatorname{ind}(u_t)\) is locally constant. Since the interval \([0,1]\) is connected, this map is constant. We therefore have \(\operatorname{ind}(u) = \operatorname{ind}(u + h)\) .
\item
  We assume that E = F. \(u = 1_{E}\) . By prop. 5 of III, p. 5, there exists a Banach space G, a continuous linear map \(p: E \to G\) and a compact linear map \(q: G \to E\) such that \(h = q \circ p\) . The endomorphism \(p \circ q\) of G is compact. By A), \(1_{G} + p \circ q\) is a Fredholm endomorphism of G with index zero. But then \(1_{E} + h = 1_{E} + q \circ p\) is a Fredholm endomorphism of E with index zero (III, p. 49, prop. 10, f)).
\item
  General case. Let \(v\) be a quasi-inverse of \(u\) . The endomorphisms \(u \circ v\) and \((u + h) \circ v\) of \(F\) differ from \(1_{F}\) by compact linear maps. By B), these are Fredholm endomorphisms of \(F\) of index zero. It follows that \(u + h\) is a Fredholm endomorphism (III, p. 40, n°2) with the same index as \(u\) (III, p. 43, no. 3, formula (2)).
\end{enumerate}

COROLLARY 1. --- Let E and F be separated locally convex spaces and \(u \in \mathcal{L}(\mathrm{E};\mathrm{F})\) . For u to be a Fredholm operator, it is necessary and sufficient that there exist a map \(v \in \mathcal{L}(\mathrm{F};\mathrm{E})\) such that the linear maps \(1_{\mathrm{E}} - v \circ u\) and \(1_{\mathrm{F}} - u \circ v\) are compact.

Since a continuous linear map of finite rank is compact, the stated condition is necessary.

Let \(v\) be an element of \(\mathcal{L}(\mathrm{F};\mathrm{E})\) such that the linear maps \(1_{\mathrm{E}} - v \circ u\) and \(1_{\mathrm{F}} - u \circ v\) are compact. By Theorem 3, \(v \circ u\) and \(u \circ v\) are Fredholm maps. Let \(p\) and \(q\) be quasi-inverses of \(v \circ u\) and \(u \circ v\) respectively. Set \(w = p \circ v\) and \(w' = v \circ q\) . Using the notation of No. 2 of III, p. 40, we have \(w \circ u \equiv 1_{\mathrm{E}}\) and \(u \circ w' \equiv 1_{\mathrm{F}}\) , whence \(w \equiv w \circ u \circ w' \equiv w'\) . It follows that \(w\) is a quasi-inverse of \(u\) .

COROLLARY 2. --- Let E and F be separated locally convex spaces and \(u \in \mathcal{L}(\mathrm{E};\mathrm{F})\) . The following conditions are equivalent:

\begin{enumerate}
\def\labelenumi{(\roman{enumi})}
\item
  The map u is a Fredholm map of index zero;
\item
  There exists an isomorphism v from E onto F such that u - v has finite rank;
\item
  There exists an isomorphism \(v\) of \(E\) onto \(F\) such that \(u - v\) is compact.
\item
  \(\Longrightarrow\) (ii): Suppose that u is a Fredholm operator of index zero. There exist topological direct sum decompositions \(E = E_{1} \oplus E_{2}\) and \(F = F_{1} \oplus F_{2}\) with \(E_{2}\) and \(F_{2}\) of finite dimension such that u vanishes on \(E_{2}\) and defines by restriction an isomorphism \(u_{1}\) of \(E_{1}\) onto \(F_{1}\) (III, p. 42, prop. 2). If \(\text{ind}(u) = 0\) then the dimensions of \(E_{2}\) and \(F_{2}\) are equal and there exists an isomorphism v from E onto F extending \(u_{1}\) . The kernel of u - v contains \(E_{1}\) , so u - v is of finite rank.
\end{enumerate}

We have (ii) \(\Longrightarrow\) (iii) since every continuous linear map of finite rank from E into F is compact, and (iii) \(\Longrightarrow\) (i) by Thm. 3.

\subsection{Perturbation of Riesz endomorphisms}

Let E be a Banach space. Recall (cf. III, p. 5, prop. 3) that the set \(\mathcal{L}^c (\mathrm{E})\) of compact endomorphisms of E forms a closed two-sided ideal of the Banach algebra \(\mathcal{L}(\mathrm{E})\) . The quotient Banach algebra is called the Calkin algebra of E; it is denoted \(\mathcal{Calk}(\mathrm{E})\) . We denote by \(\pi\) the canonical homomorphism of \(\mathcal{L}(\mathrm{E})\) onto \(\mathcal{Calk}(\mathrm{E})\) . By cor. 1 of III, p. 74, an endomorphism \(u\) of E is a Fredholm endomorphism if and only if \(\pi (u)\) is invertible in the algebra \(\mathcal{Calk}(\mathrm{E})\) .

PROPOSITION 2. --- Let \(u \in \mathcal{L}(\mathrm{E})\) such that \(\|1_{\mathrm{E}} - \pi(u)\| < 1\) in the algebra \(\mathcal{Calk}(\mathrm{E})\) . Then \(u\) is a Riesz endomorphism of E.

Let \(r \geqslant 1\) be an integer such that \(\| 1_{\mathrm{E}} - \pi(u)\|^{r} < \frac{1}{2}\) . Let \(\mathrm{P} \in \mathbf{C}[\mathrm{X}]\) be the polynomial \(\frac{1 - (1 - \mathrm{X})^r}{\mathrm{X}}\) . Set \(v = 1_{\mathrm{E}} - (1_{\mathrm{E}} - u)^r\) . We therefore have \(v = u\mathrm{P}(u)\) and \(\| 1_{\mathrm{E}} - \pi(v)\| < \frac{1}{2}\) . Since the endomorphisms \(u\) and \(\mathrm{P}(u)\) of E commute, it suffices to prove that \(v\) is a Riesz endomorphism of E (III, p. 49, prop. 9).

Since \(\|1_{\mathrm{E}}-\pi(v)\|<1/2\) , there exists, by definition of the quotient norm in the space \(\mathcal{Calk}(\mathrm{E})\) , a compact endomorphism h of E and an endomorphism w of E such that \(v=1_{E}+h+w\) and \(\|w\|<\frac{1}{2}\) . By corollary 1 of I, p. 22, the element \(1_{E}+w\) is an automorphism of E. Since h is compact, \(v=(1_{\mathrm{E}}+w)+h\) is a Fredholm endomorphism of E of index 0 (cor. 2 of III, p. 74). For every integer \(n\geqslant0\) , denote by \(N_{n}\) the kernel of \(v^{n}\) . To prove that v is a Riesz endomorphism of E, it suffices to show that there exists an integer \(n \geqslant 0\) such that \(N_{n} = N_{n+1}\) (III, p. 46, def. 2 and III, p. 46, remark).

Let us argue by contradiction by assuming the sequence \((\mathrm{N}_{n})\) strictly increasing. For every \(n \in N\) , denote by \(p_{n}\) the canonical map from E onto the normed space \(E/N_{n}\) . Let c be a real number such that \(2\|w\| < c < 1\) . Let \(n \in N\) . Since \(N_{n+1}\) is different from \(N_{n}\) , there exists an element \(x_{n} \in N_{n+1}\) such that \(\|p_{n}(x_{n})\| = c\) and such that \(\|x_{n}\| < 1\) (indeed, there exists \(y \in N_{n+1}/N_{n}\) of norm c, so for every \(\varepsilon > 0\) , there exist \(x_{n} \in N_{n+1}\) such that \(p_{n}(x_{n}) = y\) and \(\|x_{n}\| \leqslant c + \varepsilon\) ).

Let \(m\) , \(n\) be two integers such that \(m > n \geqslant 0\) . We have

\[
\begin{array}{r l} & {\| h (x _ {m}) - h (x _ {n}) \| = \| v (x _ {m} - x _ {n}) - (1 _ {\mathrm{E}} + w) (x _ {m} - x _ {n}) \|} \\ & {\qquad \geqslant \| (v - 1 _ {\mathrm{E}}) (x _ {m} - x _ {n}) \| - \| w (x _ {m} - x _ {n}) \|} \\ & {\qquad \geqslant \| (v - 1 _ {\mathrm{E}}) (x _ {m} - x _ {n}) \| - 2 \| w \|} \end{array}\tag{1}
\]

since \(\| x_m\| \leqslant 1\) and \(\| x_n\| \leqslant 1\) . Moreover, note that

\[
(v - 1 _ {\mathrm{E}}) (x _ {m} - x _ {n}) = v (x _ {m} - x _ {n}) + x _ {n} - x _ {m}.
\]

But since n \textless{} m and \(v(\mathrm{N}_{m+1}) \subset \mathrm{N}_{m}\) , we have \(v(x_{m} - x_{n}) + x_{n} \in \mathrm{N}_{m}\) , whence

\[
\left\| v (x _ {m} - x _ {n}) + x _ {n} - x _ {m} \right\| \geqslant \left\| p _ {m} (x _ {m}) \right\| = c.
\]

Inequality (1) therefore yields the lower bound

\[
\left\| h \left(x _ {m}\right) - h \left(x _ {n}\right) \right\| \geqslant c - 2 \| w \| > 0.\tag{2}
\]

The sequence \((h(x_{m}))_{m\in\mathbf{N}}\) therefore has no cluster point in E, which contradicts the fact that the image under h of the unit ball of E is relatively compact, since h is compact (remark 1 of III, p. 2).

COROLLARY 1. --- Let E be a separated locally convex space and let \(u \in \mathcal{L}(\mathrm{E})\) . For u to be a Riesz endomorphism, it is necessary and sufficient that there exist an element v of \(\mathcal{L}(\mathrm{E})\) which commutes with u and such that \(1_{\mathrm{E}} - u \circ v\) is compact.

Since every continuous linear map of finite rank is compact, the condition is necessary (III, p. 47, prop. 6 d)). Let us prove that it is sufficient. Let v be an element of \(\mathcal{L}(\mathrm{E})\) that commutes with u and such that the endomorphism \(h = 1_{E} - u \circ v\) of E is compact. There exist a Banach space G, a continuous linear map \(p: E \to G\) and a compact linear map \(q: G \to E\) such that \(h = q \circ p\) (III, p. 5, prop. 5). The endomorphism \(p \circ q\) of G is compact. By prop. 2, the linear map \(1_{G} - p \circ q\) is a Riesz endomorphism of G. But then \(u \circ v = 1_{E} - q \circ p\) is a Riesz endomorphism of E (III, p. 49, prop. 10, g)). Since u and v commute, u is a Riesz endomorphism of E (III, p. 49, prop. 9).

COROLLARY 2. --- Let E be a separated locally convex space, u a Riesz endomorphism of E, and \(h \in \mathcal{L}(\mathrm{E})\) a compact linear map. If u and h commute, then \(u + h\) is a Riesz endomorphism of E.

Let \(v\) be the canonical quasi-inverse of \(u\) . It commutes with \(u\) and with \(h\) (III, p. 47, prop. 6), hence with \(u + h\) . The endomorphism \(1_{\mathrm{E}} - (u + h) \circ v\) of E is the sum of the continuous linear map of finite rank \(1_{\mathrm{E}} - u \circ v\) and the compact linear map \(-h \circ v\) , hence is compact. It follows that \(u + h\) is a Riesz endomorphism of E (corollary 1).

PROPOSITION 3. --- Let E be a Banach space. Let u be an endomorphism of E, and let A be its commutant in \(\mathcal{L}(\mathrm{E})\) . The closure B of \(\pi(\mathrm{A})\) in the algebra \(\mathcal{Calk}(\mathrm{E})\) is a Banach algebra. For u to be a Riesz endomorphism of E, it is necessary and sufficient that \(\pi(u)\) be invertible in B.

Since \(\pi(A)\) is a unital subalgebra of \(\mathcal{Calk}(E)\) , its closure B is a Banach algebra. If \(u\) is a Riesz endomorphism of E, the element \(\pi(u)\) has an inverse in \(\pi(A)\) (cor. 1), hence in B. Conversely, suppose that \(\pi(u)\) has an inverse \(s\) in B. By definition of B, there exists an element \(v\) of A such that

\[
\left\| s - \pi (v) \right\| <   \frac {1}{\left\| \pi (u) \right\|},
\]

whence \(\|1_{E}-\pi(u\circ v)\|<1\) . By prop. 2, the linear map \(u\circ v\) is a Riesz endomorphism of E. The same is true of u since u and v commute (III, p. 49, prop. 9).

\subsection{The theory of Frédéric Riesz}

Let E be a separated locally convex space and let h be a compact endomorphism of E. Let \(\lambda \in K\) . The endomorphism \(\lambda h\) is compact, hence \(1_{E} - \lambda h\) is a Riesz endomorphism of E (cor. 2 of prop. 2 of III, p. 75); denote by \(N_{\lambda}\) and \(I_{\lambda}\) its nilspace and its conilspace. By prop. 6 of III, p. 47, we have the following properties:

\begin{enumerate}
\def\labelenumi{(\roman{enumi})}
\item
  The vector subspaces \(I_{\lambda}\) and \(N_{\lambda}\) of E are closed and stable under h;
\item
  The locally convex space E is the topological direct sum of \(I_{\lambda}\) and \(N_{\lambda}\) ;
\item
  The map \(1_{\mathrm{E}} - \lambda h\) induces by restriction an automorphism of \(\mathrm{I}_{\lambda}\) ;
\item
  The vector space \(N_{\lambda}\) is of finite dimension, and there exists an integer \(n_{\lambda} \geqslant 0\) such that \((1_{\mathrm{E}} - \lambda h)^{n_{\lambda}}(x) = 0\) for all \(x \in N_{\lambda}\) .
\end{enumerate}

These properties determine the spaces \(I_{\lambda}\) and \(N_{\lambda}\) in a unique manner (A, VIII, p. 26, remark 2).

Lemma 3. --- Let X be a topological space whose topology has a countable base. Every discrete subset of X is countable.

Let \(\mathcal{B}\) be a countable base of the topology of X and let D be a discrete subset of X. For each element \(x\) of D, there exists an element \(\mathrm{B}_x\) of \(\mathcal{B}\) such that \(\mathrm{B}_x \cap \mathrm{D} = \{x\}\) . The map from D to \(\mathcal{B}\) defined by \(x \mapsto \mathrm{B}_x\) is injective, whence the lemma.

THEOREM 4. --- Let E be a separated locally convex space and let h be a compact endomorphism of E. The set \(\Sigma\) of \(\lambda \in K\) such that \(1_{\mathrm{E}} - \lambda h\) is not an automorphism of E is a countable, closed, and discrete subset of the field K.

The set \(\Sigma\) is also the set of \(\lambda\in\mathrm{K}\) such that \(1_{\mathrm{E}}-\lambda h\) is not injective and the set of \(\lambda\in\mathrm{K}\) such that \(1_{\mathrm{E}}-\lambda h\) is not surjective.

There exists a Banach space G, a continuous linear map \(p: E \rightarrow G\) and a compact linear map \(q: G \rightarrow E\) such that \(h = q \circ p\) (III, p. 5, prop. 5). The endomorphism \(h' = p \circ q\) of G is compact, and \(1_{E} - \lambda h\) is an automorphism of E if and only if \(1_{G} - \lambda h'\) is an automorphism of G (prop. 10 of III, p. 49, e)). It therefore suffices for us to prove the theorem when E is a Banach space, which we henceforth assume.

Let \(\lambda \in K\) . Since the index of a Riesz endomorphism is zero, it is equivalent to say that the map \(1_{\mathrm{E}} - \lambda h\) is an automorphism of E, or that it is injective, or again that it is surjective.

The set \(\Sigma\) is closed in K since the set of automorphisms of E is open in \(\mathcal{L}(\mathrm{E})\) . Let \(\lambda\) be an element of \(\Sigma\) . We have \(\lambda \neq 0\) . Let N denote the nilspace of h and I its conilspace, and denote by \(h_{I}\) and \(h_{N}\) the endomorphisms of I and of N induced by h on these subspaces. Then \(1_{I} - \lambda h_{I}\) is an automorphism of I, and there exists a neighborhood U of \(\lambda\) in K such that \(1_{I} - \mu h_{I}\) is an automorphism of I for all \(\mu \in U\) . The endomorphism \(1_{N} - \lambda h_{N}\) of N is nilpotent; for all \(\mu \neq \lambda\) , we have

\[
1 _ {\mathrm{N}} - \mu h _ {\mathrm{N}} = \frac {\lambda - \mu}{\lambda} \Big (1 _ {\mathrm{N}} + \frac {\mu}{\lambda - \mu} (1 _ {\mathrm{N}} - \lambda h _ {\mathrm{N}}) \Big),
\]

hence \(1_{N} - \mu h_{N}\) is an automorphism of N. Consequently, the set \(U \cap \Sigma\) is reduced to the single element \(\lambda\) and the set \(\Sigma\) is discrete. It is countable by Lemma 3.

\subsection{Fredholm alternative}

PROPOSITION 4 (Fredholm alternative). --- Let E be a Banach space and let h be a compact endomorphism of E. Let \(\lambda\) be an element of K-\{0\} . Let F be the kernel of \(1_{\mathrm{E}} - \lambda h\) and let G be the kernel of \(1_{\mathrm{E}'} - \lambda^{t}h\) .

\begin{enumerate}
\def\labelenumi{\alph{enumi})}
\item
  The spaces F and G have the same finite dimension n ≥ 0;
\item
  For \(y \in \mathrm{E}\) , there exists \(x \in \mathrm{E}\) such that \(x - \lambda h(x) = y\) if and only if \(\langle y, \ell \rangle = 0\) for all \(\ell \in \mathrm{G}\) ;
\item
  For \(\ell \in \mathrm{E}'\) , there exists \(f \in \mathrm{E}'\) such that \(f - \lambda^t h(f) = \ell\) if and only if \(\langle x, f \rangle = 0\) for all \(x \in \mathrm{F}\) .
\end{enumerate}

In particular, one and only one of the following conditions holds:

\begin{enumerate}
\def\labelenumi{(\roman{enumi})}
\item
  The space F is nonzero;
\item
  For every \(y \in \mathrm{E}\) , there exists \(x \in \mathrm{E}\) such that \(x - \lambda h(x) = y\) .
\end{enumerate}

By replacing h with \(\lambda h\) we reduce to the case where \(\lambda = 1\) .

The endomorphism \(^{t}h\) is compact by Corollary 1 of III, p. 9. The endomorphisms \(w = 1_{\mathrm{E}} - h\) of E and \(w' = 1_{\mathrm{E}'} - ^{t}h\) of \(\mathrm{E}'\) are therefore Riesz endomorphisms (cor. 2 of prop. 2 of III, p. 75). In particular, their kernels F and G are finite-dimensional. Since \(w' = ^{t}w\) and since the index of w is zero, the dimension of F is equal to that of G by formula (3) of III, p. 43.

For an element \(y\) of E, the equation \(x - h(x) = y\) has a solution \(x \in \mathrm{E}\) if and only if \(x\) belongs to the image of \(w\) . Since the dual of the cokernel of \(w\) is identified with the kernel G of \(\operatorname{Ker}(^t w)\) (EVT, IV, p. 27, prop. 2), this is the case if and only if \(\langle y, \ell \rangle = 0\) for all \(\ell \in \mathrm{G}\) .

For an element \(\ell\) of E', the equation \(f - {}^t h(f) = \ell\) has a solution \(f \in \mathrm{E}'\) if and only if \(f\) belongs to the image of \(w'\) . Since the image of \({}^t w\) is the orthogonal of the kernel F of \(w\) (EVT, IV, p. 27, prop. 2) , this is the case if and only if \(\langle x, f \rangle = 0\) for all \(x \in \mathrm{F}\) .

Example. --- Let X be a compact space, \(\mu\) a measure on X, real or complex depending on whether K is equal to R or C, and N: \(X \times X \to K\) a continuous function.

Let E denote the Banach space \(\mathcal{C}(\mathrm{X};\mathrm{K})\) . Suppose that a function \(a\in\mathrm{E}\) and \(\lambda\in\mathrm{K}-\{0\}\) are given. We study the problem of the existence and uniqueness of solutions \(f\in\mathrm{E}\) of the integral equation

\[
f (x) - \lambda \int_ {\mathrm{X}} \mathrm{N} (x, y) f (y) d \mu (y) = a (x) \quad (x \in \mathrm{X}).\tag{3}
\]

For this, let us introduce the set \(F_{\lambda}\) of solutions \(f \in E\) of the associated homogeneous equation

\[
f (x) - \lambda \int_ {\mathrm{X}} \mathrm{N} (x, y) f (y) d \mu (y) = 0 \qquad (x \in \mathrm{X}),\tag{4}
\]

and the set \(G_{\lambda}\) of solutions \(g \in E\) of the homogeneous transposed equation of (4), that is,

\[
g (y) - \lambda \int_ {\mathrm{X}} \mathrm{N} (x, y) g (x) d \mu (x) = 0 \qquad (y \in \mathrm{X}).\tag{5}
\]

If it is not empty, the set of solutions of (3) is an affine subspace of E with direction \(F_{\lambda}\) .

For \(f\in \mathrm{E}\) and \(x\in \mathrm{X}\) , set

\[
h (f) (x) = \int_ {\mathrm{X}} \mathrm{N} (x, y) f (y) d \mu (y);
\]

the map \(f \mapsto h(f)\) thus defined is a compact endomorphism of the Banach space E (cor. of prop. 2 of III, p. 26). The dual E' of E consists of measures on X, real or complex depending on whether K equals R or C (INT, III, p. 47, § 1, n° 3, def. 2). The transpose \(^{t}h\) of h is the endomorphism of E' characterized by the relation \(\langle f, ^{t}h(m)\rangle = \langle h(f), m\rangle\) for \(f \in E\) and \(m \in E'\) , that is,

\[
\begin{array}{r l} \int_ {\mathrm{X}} f (x) d (^ {t} h (m)) (x) & = \int_ {\mathrm{X}} h (f) (x) d m (x) \\ & = \int_ {\mathrm{X}} \Big (\int_ {\mathrm{X}} \mathrm{N} (x, y) f (y) d \mu (y) \Big) d m (x) \\ & = \int_ {\mathrm{X}} \Big (\int_ {\mathrm{X}} \mathrm{N} (x, y) d m (x) \Big) f (y) d \mu (y) \end{array}\tag{6}
\]

for \(m \in E'\) and \(f \in E\) (INT, III, p. 84, § 4, no. 1, th. 2).

Lemma 4. --- The linear map from \(G_{\lambda}\) into \(E'\) defined by \(g \mapsto g \cdot \mu\) is injective, and its image is the kernel of \(1_{E'} - \lambda^{t}h\) .

If \(g \in G_{\lambda}\) satisfies \(g \cdot \mu = 0\) then formula (5) implies \(g = 0\) , so the map \(g \mapsto g \cdot \mu\) is injective.

Let \(m \in \mathrm{E}'\) be a measure belonging to the kernel of \(1_{\mathrm{E}'} - \lambda^t h\) . By relation (6), we have

\[
\int_ {\mathrm{X}} f (x) d m (x) = \lambda \int_ {\mathrm{X}} \left(\int_ {\mathrm{X}} \mathrm{N} (x, y) d m (x)\right) f (y) d \mu (y)\tag{7}
\]

for every continuous function \(f: \mathrm{X} \to \mathrm{K}\) . The measure \(m\) is therefore equal to the measure \(g \cdot \mu\) , where \(g\) is the continuous function from \(\mathrm{X}\) into \(\mathrm{K}\) defined by

\[
g (y) = \lambda \int_ {\mathrm{X}} \mathrm{N} (x, y) d m (x)\tag{8}
\]

for all \(y \in X\) . Formula (8) then implies that the function g belongs to \(G_{\lambda}\) .

Conversely, let \(g: X \to K\) be a continuous function belonging to \(G_{\lambda}\) and set \(m = g \cdot \mu\) . For every continuous function \(f: X \to K\) , we therefore have

\[
\begin{array}{l} \int_ {\mathrm{X}} f (y) d m (y) = \int_ {\mathrm{X}} g (y) f (y) d \mu (y) \\ \qquad = \lambda \int_ {\mathrm{X}} \Big (\int_ {\mathrm{X}} \mathrm{N} (x, y) g (x) d \mu (x) \Big) f (y) d \mu (y) \\ \qquad = \lambda \int_ {\mathrm{X}} f (x) d m (x), \end{array}
\]

so that \({}^{t}h(m)=\lambda m\) . This proves Lemma 4.

Applying Proposition 4 and Theorem 4 of III, p. 78 to the endomorphism h, we then obtain the following statements.

THEOREM 5. --- Let \(\lambda\in K\) .

\begin{enumerate}
\def\labelenumi{\alph{enumi})}
\item
  The sets \(F_{\lambda}\) and \(G_{\lambda}\) are finite-dimensional vector subspaces of \(\mathcal{C}(X;K)\) and their dimensions are equal;
\item
  For equation (3) to have at least one solution \(f \in \mathcal{C}(\mathrm{X}, \mathrm{K})\) it is necessary and sufficient that one has \(\int_{\mathrm{X}} a(x) g(x) d\mu(x) = 0\) for every function \(g \in G_{\lambda}\) . The set of solutions of (3) is then an affine subspace of \(\mathcal{C}(\mathrm{X}; \mathrm{K})\) with direction \(F_{\lambda}\) ;
\item
  One of the following mutually exclusive conditions is satisfied:
\end{enumerate}

\begin{enumerate}
\def\labelenumi{(\roman{enumi})}
\item
  For any function \(a \in \mathcal{C}(\mathrm{X};\mathrm{K})\) there exists a unique solution \(f \in \mathcal{C}(\mathrm{X};\mathrm{K})\) of the integral equation (3);
\item
  The homogeneous equation (4) has a nonzero solution \(f \in \mathcal{C}(\mathrm{X};\mathrm{K})\) .
\end{enumerate}

Corollary. --- The set of \(\lambda\in\mathrm{K}\) for which the space \(\mathrm{F}_{\lambda}\) is nonzero is a countable, closed, and discrete subset of K.

\section{SPECTRAL PROPERTIES OF ENDOMORPHISMS OF BANACH SPACES}

Unless stated otherwise, the vector spaces considered in this paragraph are vector spaces over C. The spectrum of an endomorphism of a complete normable space E is the spectrum relative to the unital algebra \(\mathcal{L}(\mathrm{E})\) (cf. no. 7 of I, p. 127).

\subsection{Isolated points and sensitive points of the spectrum}

Let E be a complete normable space and let u be an endomorphism of E. Let \(\lambda\) be an isolated point of the spectrum of u. Recall that we denote by \(e_{\lambda}(u)\) the spectral projector associated with u and the closed and open set \(\{\lambda\}\) of the spectrum of u (no. 3 of I, p. 131). The endomorphism \(u - \lambda1_{E}\) induces an automorphism of \(\operatorname{Ker}(e_{\lambda}(u))\) and a quasi-nilpotent endomorphism of \(\operatorname{Im}(e_{\lambda}(u))\) (loc. cit.).

DEFINITION 1. --- One says that \(\lambda\) has finite spectral multiplicity for \(u\) if the spectral projector \(e_{\lambda}(u)\) is of finite rank; in this case, the integer \(m_{\lambda}(u) = \dim (\mathrm{Im}(e_{\lambda}(u)))\) is called the spectral multiplicity of \(\lambda\) for \(u\) .

DEFINITION 2. --- One calls sensitive points of the spectrum of u the isolated points of Sp(u) of finite spectral multiplicity. The set of sensitive points of Sp(u) is called the sensitive spectrum of u and is denoted Sp \(_{s}\) (u). \(^{(2)}\) One calls essential spectrum of u and denotes by Sp \(_{e}\) (u) the complementary set \(\mathrm{Sp}(u) - \mathrm{Sp}_{\mathrm{s}}(u)\).

As the set \(\mathrm{Sp}_{\mathrm{s}}(u)\) consists of isolated points of \(\mathrm{Sp}(u)\) , it is open in \(\mathrm{Sp}(u)\) , discrete, and countable (III, p. 78, lemma 3); it is also bounded. The set \(\mathrm{Sp}_{\mathrm{e}}(u)\) is closed in \(\mathrm{Sp}(u)\) , hence compact.

By prop. 14 of III, p. 54, the sensitive points of the spectrum of \(u\) are the complex numbers \(\lambda\) such that \(u - \lambda 1_{\mathrm{E}}\) is a Riesz endomorphism of E that is not an automorphism of E.

For every complex number \(\lambda\) , denote by \(\mathrm{N}_{\lambda}(u)\) and \(\mathrm{I}_{\lambda}(u)\) the nilspace and the conilspace of the endomorphism \(u - \lambda 1_{\mathrm{E}}\) of E. Since E is a Fréchet space, \(u - \lambda 1_{\mathrm{E}}\) is a Riesz endomorphism if and only if \(\mathrm{N}_{\lambda}(u)\) is finite-dimensional and \(\mathrm{I}_{\lambda}(u)\) has finite codimension (prop. 13 of III, p. 53); when this is the case, \(u - \lambda 1_{\mathrm{E}}\) is invertible if and only if \(\mathrm{N}_{\lambda}(u)\) is reduced to 0 (III, p. 47, prop. 6). The sensitive points of the spectrum of \(u\) are therefore the complex numbers \(\lambda\) such that \(\mathrm{N}_{\lambda}(u)\) is of nonzero finite dimension and \(\mathrm{I}_{\lambda}(u)\) is of finite codimension in E (III, p. 54, prop. 14). The spectral multiplicity \(m_{\lambda}(u)\) of \(\lambda\) is then the dimension of \(\mathrm{N}_{\lambda}(u)\) , and the endomorphism \(u - \lambda 1_{\mathrm{E}}\) induces by restriction an automorphism of \(\mathrm{I}_{\lambda}(u)\) and a nilpotent endomorphism of \(\mathrm{N}_{\lambda}(u)\) . Since \(\mathrm{N}_{\lambda}(u)\) is not reduced to 0, it follows that \(\lambda\) is an eigenvalue of \(u\) . The smallest integer \(n \geqslant 1\) such that \(\operatorname{Ker}((u - \lambda 1_{\mathrm{E}})^n)\) is equal to \(\mathrm{N}_{\lambda}(u)\) is the order of the pole of the resolvent of \(u\) at the point \(\lambda\) (No. 3 of I, p. 131), which is bounded above by \(m_{\lambda}(u)\) .

In particular, we have therefore proved:

PROPOSITION 1. --- The sensitive points of the spectrum of u are eigenvalues of finite spectral multiplicity.

The essential spectrum of \(u\) consists of the complex numbers \(\lambda\) such that \(u - \lambda 1_{\mathrm{E}}\) is not a Riesz endomorphism of E. In particular, if \(u\) is compact, then every \(\lambda \in \mathrm{Sp}(u) - \{0\}\) is a sensitive point of the spectrum (cor. 2 of III, p. 77).

PROPOSITION 2. --- Let E be a complete normable space and u an endomorphism of E. Let H be a finite subset of \(\mathrm{Sp}_{\mathrm{s}}(u)\) . The set H is open and closed in \(\mathrm{Sp}(u)\) . The spectral projector \(e_{\mathrm{H}}(u)\) has finite rank, has image \(\sum_{\lambda \in \mathrm{H}}\mathrm{N}_{\lambda}(u)\) and kernel \(\bigcap_{\lambda \in \mathrm{H}}\mathrm{I}_{\lambda}(u)\) .

This follows from no. 3 of I, p. 131 since, for every \(\lambda \in \mathrm{Sp}_{\mathrm{s}}(u)\) , the subspaces \(\mathrm{N}_{\lambda}(u)\) and \(\mathrm{I}_{\lambda}(u)\) are respectively the image and the kernel of the spectral projector \(e_{\lambda}(u)\) .

COROLLARY 1. --- Let V be a neighborhood of \(\mathrm{Sp_e}(u)\) .

\begin{enumerate}
\def\labelenumi{\alph{enumi})}
\item
  There exists a decomposition of E as a topological direct sum F ⊕ G such that F is finite-dimensional, F and G are stable under every endomorphism commuting with u, and the spectrum of the endomorphism of G induced by u is contained in V.
\item
  Suppose V is nonempty. There exists an element v of the bicommutant of u in \(\mathcal{L}(\mathrm{E})\) whose spectrum is contained in V and such that \(v - u\) has finite rank.
\end{enumerate}

Set \(\mathrm{H} = \mathrm{Sp}(u) \cap (\mathbf{C} - \mathring{\mathrm{V}})\) . The set H is contained in \(\mathrm{Sp}_{\mathrm{s}}(u)\) , hence is discrete; since it is closed in the compact space \(\mathrm{Sp}(u)\) , it is finite. One may take for F and G the image and the kernel of the spectral projector \(e_{\mathrm{H}}(u)\) (prop. 2).

Suppose that V is not empty. Let \(\mu\in\mathrm{V}\) and denote by \(v\) the endomorphism of E that coincides with the homothety of ratio \(\mu\) in F and with \(u\) in G. Its spectrum is contained in \(\{\mu\}\cup(\mathrm{Sp}(u)-\mathrm{H})\) , therefore in V, and \(v-u\) is of finite rank since F is finite-dimensional.

For every compact subset S of C, we denote by \(\mathcal{O}(S)\) the algebra of germs of holomorphic functions in a neighborhood of S and with values in C (I, p. 49, §4, no. 1).

COROLLARY 2. --- Let \(f \in \mathcal{O}(\mathrm{Sp}(u))\) and let \(\mu\) be a complex number. Let H be the set of \(\lambda \in \mathrm{Sp}(u)\) such that \(f(\lambda) = \mu\) . Then \(\mu\) is a sensitive point of the spectrum of \(f(u)\) if and only if H is finite, nonempty, and contained in the sensitive spectrum of \(u\) . Under these conditions, the spectral projector \(e_{\mu}(f(u))\) is equal to the projector \(e_{\mathrm{H}}(u)\) associated with \(u\) and H. We have

\[
\mathrm{N} _ {\mu} (f (u)) = \bigoplus_ {\lambda \in \mathrm{H}} \mathrm{N} _ {\lambda} (u), \quad \mathrm{I} _ {\mu} (f (u)) = \bigcap_ {\lambda \in \mathrm{H}} \mathrm{I} _ {\lambda} (u)
\]

and the spectral multiplicity of \(\mu\) for \(f(u)\) is the sum of the spectral multiplicities of the elements of H, that is,

\[
m _ {\mu} (f (u)) = \sum_ {\lambda \in \mathrm{H}} m _ {\lambda} (u).
\]

The complex number \(\mu\) belongs to \(f(\mathrm{Sp}(u))\) , and therefore to the spectrum of \(f(u)\) (prop. 8 of I, p. 75), if and only if H is nonempty. We then have \(e_{\mu}(f(u)) = e_{\mathrm{H}}(u)\) (no. 1 of I, p. 127). The other assertions follow from this by Prop. 2.

Examples. --- 1) Let E be a finite-dimensional vector space and u an endomorphism of E. The spectrum of u is finite and coincides with \(\mathrm{Sp}_{\mathrm{s}}(u)\) . Its elements are the roots of the characteristic polynomial \(\chi_{u}\) of u; these are the eigenvalues of u. The spectral multiplicity \(m_{\lambda}(u)\) of an element \(\lambda \in \mathrm{Sp}(u)\) is the multiplicity of \(\lambda\) as a root of the polynomial \(\chi_{u}\) (A, VII, p. 36, cor.). By cor. 2 above, we have

\[
\chi_ {f (u)} (\mathrm{T}) = \prod_ {\lambda \in \operatorname{Sp} (u)} (\mathrm{T} - f (\lambda)) ^ {m _ {\lambda}}
\]

for all \(f \in \mathcal{O}(\mathrm{Sp}(u))\) . This generalizes prop. 10 of A, VII, p. 36.

\begin{enumerate}
\def\labelenumi{\arabic{enumi})}
\setcounter{enumi}{1}
\tightlist
\item
  Let X be a compact subset of C and \(\mathcal{C}(X)\) the Banach space of continuous functions on X with complex values. Denote by u the endomorphism of \(\mathcal{C}(X)\) that sends \(f \in \mathcal{C}(X)\) to the function \(z \mapsto zf(z)\) in \(\mathcal{C}(X)\) . The spectrum of u is equal to X, its sensitive points are the isolated points of X and their spectral multiplicities are equal to 1.
\end{enumerate}

\subsection{A partition of the spectrum}

Denote by \(\overline{Z}\) the subset \(Z \cup \{-\infty, +\infty\}\) of \(\overline{R}\) . If u is a linear map whose kernel or cokernel is finite-dimensional, one calls the index of u the element \(\text{ind}(u)\) of \(\overline{Z}\) defined by

\[
\operatorname{ind} (u) = \dim \operatorname{Coker} (u) - \dim \operatorname{Ker} (u)
\]

(cf. no. 6 of III, p. 67).

DEFINITION 3. --- Let E be a complete normable space and u an endomorphism of E. For every \(n \in \overline{Z}\) , denote by \(\mathrm{Sp}_{n}(u)\) the set of complex numbers \(\lambda \in \mathrm{Sp}_{\mathrm{e}}(u)\) such that \(u - \lambda1_{E}\) is a strict morphism, whose kernel or cokernel is finite-dimensional, and whose index is \(n\). We denote \(\mathrm{Sp}_{\omega}(u)\) the complement in \(\mathrm{Sp}_{\mathrm{e}}(u)\) of the union of the subsets \(\mathrm{Sp}_{n}(u)\) for \(n \in \overline{Z}\) .

The sets \(\mathrm{Sp}_{\mathrm{s}}(u)\) , \(\mathrm{Sp}_{n}(u)\) for \(n \in \overline{Z}\) and \(\mathrm{Sp}_{\omega}(u)\) form a partition of the spectrum of u.

Every endomorphism of E whose cokernel is finite-dimensional is strict (III, p. 52, lemma 6). The Fredholm endomorphisms of E are the endomorphisms of E whose kernel and cokernel are finite-dimensional (III, p. 52, prop. 11). The set \(\mathbf{C} - \mathrm{Sp}_{\mathrm{e}}(u)\) consists of the \(\lambda \in \mathbf{C}\) such that \(u - \lambda 1_{\mathrm{E}}\) is a Riesz endomorphism of E, and such an endomorphism is a Fredholm endomorphism of E of index 0.

Consequently, for \(\lambda \in \mathbf{C}\) and for \(n\in \mathbf{Z} - \{0\}\) , we have:

\(\lambda \in \mathbf{C} - \mathrm{Sp}(u) \iff u - \lambda 1_{\mathrm{E}}\) is an automorphism;

\(\lambda \in \mathrm{Sp}_{\mathrm{s}}(u) \quad \Longleftrightarrow \quad u - \lambda 1_{\mathrm{E}} \text{ is a Riesz endomorphism, but is not an automorphism;}\)

\(\lambda \in \mathrm{Sp}_{0}(u) \quad \Longleftrightarrow \quad u - \lambda 1_{\mathrm{E}}\) is a Fredholm endomorphism of index 0 of E but is not a Riesz endomorphism of E;

\(\lambda \in \mathrm{Sp}_n(u) \quad \Longleftrightarrow \quad u - \lambda 1_{\mathrm{E}}\) is a Fredholm endomorphism of index \(n\) of E (\(n \neq 0\));

\[
\lambda \in \mathrm{Sp} _ {- \infty} (u) \quad \Longleftrightarrow \quad u - \lambda 1 _ {\mathrm{E}} \text{ is strict, its kernel is infinite-dimensional and its cokernel is finite-dimensional;}
\]

\(\lambda \in \mathrm{Sp}_{+\infty}(u) \quad \Longleftrightarrow \quad u - \lambda 1_{\mathrm{E}}\) is strict, its kernel is finite-dimensional and its cokernel is infinite-dimensional;

\(\lambda \in \mathrm{Sp}_{\omega}(u) \quad \Longleftrightarrow \quad \text{either } u - \lambda 1_{\mathrm{E}} \text{ is not strict, or its kernel and cokernel are infinite-dimensional.}\)

Remark. --- Let \(\pi\) denote the canonical homomorphism of \(\mathcal{L}(\mathrm{E})\) onto the Calkin algebra \(\mathcal{C}alk(\mathrm{E})\) of E (cf. no. 4 of III, p. 75). The spectrum of \(\pi(u)\) with respect to the algebra \(\mathcal{C}alk(\mathrm{E})\) is sometimes called the stable spectrum of \(u\) . Its elements are the complex numbers \(\lambda\) such that \(u - \lambda 1_{\mathrm{E}}\) is not a Fredholm endomorphism of E (cor. 1 of th. 3 of III, p. 73). It is therefore equal to \(\mathrm{Sp}_{\omega}(u) \cup \mathrm{Sp}_{-\infty}(u) \cup \mathrm{Sp}_{+\infty}(u)\) .

Let A be the set of endomorphisms of E commuting with u. The closure B of \(\pi(\mathrm{A})\) in \(\mathcal{C}alk(\mathrm{E})\) is a complete normable algebra and the spectrum of \(\pi(u)\) with respect to B is the essential spectrum \(\mathrm{Sp}_{\mathrm{e}}(u)\) of u (III, p. 77, prop. 3). By the cor. of prop. 6 of I, p. 28, applied to the subalgebra B of \(\mathcal{C}alk(\mathrm{E})\) , the set \(\mathrm{Sp}_{\mathrm{e}}(u)\) is the union of \(\mathrm{Sp}_{\omega}(u) \cup \mathrm{Sp}_{-\infty}(u) \cup \mathrm{Sp}_{+\infty}(u)\) and of certain bounded connected components of its complement.

THEOREM 1. --- Let E be a complete normable space and u an endomorphism of E.

\begin{enumerate}
\def\labelenumi{\alph{enumi})}
\item
  The set \(\mathrm{Sp}_{\omega}(u)\) is compact. It is nonempty if E is of infinite dimension;
\item
  Let \(n \in \overline{Z}\) . The set \(\mathrm{Sp}_{n}(u)\) is a union of a family of bounded connected components of \(\mathbf{C} - \mathrm{Sp}_{\omega}(u)\) . It is open in C, and its boundary in C is contained in \(\mathrm{Sp}_{\omega}(u)\) .
\end{enumerate}

The set \(\mathbf{C} - \mathrm{Sp}_{\omega}(u)\) consists of the complex numbers \(\lambda \in \mathbf{C}\) such that \(u - \lambda 1_{\mathrm{E}}\) is a strict morphism whose kernel or cokernel is finite-dimensional. By Propositions 11 of III, p. 67 and 13 of III, p. 70, it is open. The set \(\mathrm{Sp}_{\omega}(u)\) is therefore closed. Since it is bounded, it is compact.

Let us prove \(b\) ). Let \(n \in \overline{\mathbf{Z}}\) . The set \(\mathrm{Sp}_n(u)\) is contained in \(\mathbf{C} - \mathrm{Sp}_{\omega}(u)\) . Let U be a connected component of \(\mathbf{C} - \mathrm{Sp}_{\omega}(u)\) which intersects \(\mathrm{Sp}_n(u)\) . The map \(\lambda \mapsto \operatorname{ind}(u - \lambda 1_{\mathrm{E}})\) of \(\mathbf{C} - \mathrm{Sp}_{\omega}(u)\) in \(\overline{\mathbf{Z}}\) being locally constant (cor. 1 of prop. 12 of III, p. 68 and cor. 1 of prop. 13 of III, p. 70), the index of \(u - \lambda 1_{\mathrm{E}}\) is equal to \(n\) for all \(\lambda \in U\) . If \(n \neq 0\) , this implies that U is contained in \(\mathrm{Sp}_n(u)\) . If \(n = 0\) , note that the set U is a connected component of \(\mathbf{C} - (\mathrm{Sp}_{\omega}(u) \cup \mathrm{Sp}_{-\infty}(u) \cup \mathrm{Sp}_{+\infty}(u))\) . Since U meets \(\mathrm{Sp}_0(u)\) and therefore \(\mathrm{Sp}_{\mathrm{e}}(u)\) , it follows from Remark 2 that the set U is contained in \(\mathrm{Sp}_{\mathrm{e}}(u)\) , and hence in \(\mathrm{Sp}_0(u)\) . In all cases, we conclude that \(\mathrm{Sp}_n(u)\) is the union of the connected components of \(\mathbf{C} - \mathrm{Sp}_{\omega}(u)\) that meet \(\mathrm{Sp}_n(u)\) . These are necessarily bounded since the set \(\mathrm{Sp}(u)\) is bounded. Consequently, \(\mathrm{Sp}_n(u)\) is open in \(\mathbf{C}\) and its boundary is contained in \(\mathrm{Sp}_{\omega}(u)\) . This proves \(b\) ).

Suppose finally that the set \(\mathrm{Sp}_{\omega}(u)\) is empty. By b), each of the sets \(\mathrm{Sp}_{n}(u)\) , for \(n \in \overline{Z}\) , is then empty. We therefore have \(\mathrm{Sp}(u) = \mathrm{Sp}_{\mathrm{s}}(u)\) . The spectrum of u is consequently discrete and compact, hence finite, and since all its points have finite spectral multiplicity, the vector space E is finite-dimensional (III, p. 83, prop. 2). This completes the proof of a).

COROLLARY. --- a) Let \(\Omega\) be the unbounded connected component of \(\mathbf{C} - \mathrm{Sp}_{\omega}(u)\) . We have \(\Omega \cap \mathrm{Sp}(u) \subset \mathrm{Sp}_{\mathrm{s}}(u)\) ;

\begin{enumerate}
\def\labelenumi{\alph{enumi})}
\setcounter{enumi}{1}
\tightlist
\item
  Every point adherent to \(\mathrm{Sp}_{\mathrm{s}}(u)\) that does not belong to \(\mathrm{Sp}_{\mathrm{s}}(u)\) belongs to \(\mathrm{Sp}_{\omega}(u)\) .
\end{enumerate}

The assertion \(a\) ) is a direct consequence of assertion \(b\) ) of th. 1. Let \(\lambda\) be a point adherent to \(\mathrm{Sp}_{\mathrm{s}}(u)\) that does not belong to \(\mathrm{Sp}_{\mathrm{s}}(u)\) . It belongs to the spectrum of \(u\) , since the latter is closed. It does not belong to any of the sets \(\mathrm{Sp}_{n}(u)\) , for \(n \in \overline{\mathbf{Z}}\) , since these are open (loc. cit.) and disjoint from \(\mathrm{Sp}_{\mathrm{s}}(u)\) . We therefore have \(\lambda \in \mathrm{Sp}_{\omega}(u)\) , whence \(b\) ).

PROPOSITION 3. --- Let E and F be complete normable spaces, \(u\colon E\to F\) and \(v\colon F\to E\) continuous linear maps.

\begin{enumerate}
\def\labelenumi{\alph{enumi})}
\item
  The traces on \(\mathbf{C} - \{0\}\) of the sets \(\mathrm{Sp}(v \circ u)\) and \(\mathrm{Sp}(u \circ v)\) (resp. \(\mathrm{Sp}_{\mathrm{s}}(v \circ u)\) and \(\mathrm{Sp}_{\mathrm{s}}(u \circ v)\) , resp. \(\mathrm{Sp}_{n}(v \circ u)\) and \(\mathrm{Sp}_{n}(u \circ v)\) for \(n \in \overline{\mathbf{Z}}\) , resp. \(\mathrm{Sp}_{\omega}(v \circ u)\) and \(\mathrm{Sp}_{\omega}(u \circ v)\) ) are equal;
\item
  Let \(\lambda\) be an element of \(\mathrm{Sp}_{\mathrm{s}}(v \circ u)\) distinct from 0. The spectral multiplicities of \(\lambda\) for \(v \circ u\) and for \(u \circ v\) are equal.
\end{enumerate}

Let \(\mu\) be a nonzero complex number and let \(n \in \overline{\mathbf{Z}}\) . For \(\mu 1_{\mathrm{E}} - v \circ u\) to be an automorphism (resp. a Riesz endomorphism, resp. a strict morphism whose kernel or cokernel is finite-dimensional and whose index is \(n\) ), it is necessary and sufficient that \(\mu 1_{\mathrm{F}} - u \circ v\) be one (III, p. 49, prop. 10). The assertion \(a\) ) then follows from the definitions.

Let \(\lambda\) be a point of \(\mathrm{Sp}_{\mathrm{s}}(v \circ u)\) distinct from 0. We have

\[
\dim \operatorname{Ker} ((\lambda 1 _ {\mathrm{E}} - v \circ u) ^ {n}) = \dim \operatorname{Ker} ((\lambda 1 _ {\mathrm{F}} - u \circ v) ^ {n})
\]

for all \(n \geqslant 0\) (loc. cit.), hence the spectral multiplicities of \(\lambda\) for \(v \circ u\) and \(u \circ v\) are equal.

\subsection{Spectrum of the transpose of an endomorphism}

PROPOSITION 4. --- Let E be a complete normable space, E' the dual space of E, and u an endomorphism of E.

\begin{enumerate}
\def\labelenumi{\alph{enumi})}
\item
  One has \(\mathrm{Sp}_{\mathrm{s}}(u) = \mathrm{Sp}_{\mathrm{s}}(^{t}u)\) , \(\mathrm{Sp}_n(u) = \mathrm{Sp}_{-n}(^t u)\) for all \(n \in \overline{\mathbf{Z}}\) and \(\mathrm{Sp}_{\omega}(u) = \mathrm{Sp}_{\omega}(^t u)\) ;
\item
  Every point of \(\mathrm{Sp}_{\mathrm{s}}(u)\) has the same spectral multiplicity for u and \({}^{t}u\) .
\end{enumerate}

The assertion \(a\) ) of prop. 3 of I, p. 131 shows that \(\mathrm{Sp}(u) = \mathrm{Sp}({}^{t}u)\) , and assertion \(c\) ) of loc. cit. implies that \(\mathrm{Sp}_{\mathrm{s}}(u) = \mathrm{Sp}_{\mathrm{s}}({}^{t}u)\) and that the assertion \(b\) ) is valid.

For every \(\lambda \in \mathbf{C}\) , Lemma 4 of III, p. 69 implies that \(u - \lambda 1_{\mathrm{E}}\) is a strict morphism if and only if \({}^{t}u - \lambda 1_{\mathrm{E}'}\) is one, and that

\[
\dim \operatorname{Coker} (^ {t} (u - \lambda 1 _ {\mathrm{E}})) = \dim \operatorname{Ker} (u - \lambda 1 _ {\mathrm{E}})
\]

\[
\dim \operatorname{Ker} (^ {t} (u - \lambda 1 _ {\mathrm{E}})) = \dim \operatorname{Coker} (u - \lambda 1 _ {\mathrm{E}})
\]

in \(\overline{Z}\) . Assertion a) follows from this, taking into account the definitions of the various parts of the spectrum (def. 3 of III, p. 85).

\subsection{Perturbation by a compact operator}

THEOREM 2. --- Let E be a complete normable space, u an endomorphism of E, and h a compact endomorphism of E. Then \(\mathrm{Sp}_{\omega}(u + h) = \mathrm{Sp}_{\omega}(u)\) and \(\mathrm{Sp}_n(u + h) = \mathrm{Sp}_n(u)\) for all \(n \in \overline{\mathbf{Z}} - \{0\}\) .

Let \(\lambda \in \mathbf{C}\) . For \(u + h - \lambda 1_{\mathrm{E}}\) to be a strict morphism whose kernel (resp. cokernel) is finite-dimensional, it is necessary and sufficient that \(u - \lambda 1_{\mathrm{E}}\) have the same property, by Theorem 1 of III, p. 72 (resp. Theorem 2 of III, p. 73). The equalities

\[
\begin{array}{c} \operatorname{Sp} _ {- \infty} (u + h) = \operatorname{Sp} _ {- \infty} (u), \quad \operatorname{Sp} _ {+ \infty} (u + h) = \operatorname{Sp} _ {+ \infty} (u), \\ \operatorname{Sp} _ {\omega} (u + h) = \operatorname{Sp} _ {\omega} (u) \end{array}
\]

result from this. Moreover, we have \(\operatorname{Sp}_{n}(u+h)=\operatorname{Sp}_{n}(u)\) by th. 3 of III, p. 73, for every \(n\in \mathbf{Z}-\{0\}\) .

COROLLARY. --- Suppose that the unbounded connected component of the complement of \(\mathrm{Sp}_{\omega}(u)\) contains 0. Then \(u + h\) is a Riesz endomorphism of E.

One has \(\mathrm{Sp}_{\omega}(u+h)=\mathrm{Sp}_{\omega}(u)\) (by th. 2), hence 0 belongs to the unbounded connected component of \(\mathbf{C}-\mathrm{Sp}_{\omega}(u+h)\) . By the cor. of th. 1 of III, p. 87, either 0 does not belong to the spectrum of \(u+h\) , or it is a sensitive point of this spectrum. In either case, \(u+h\) is a Riesz endomorphism of E.

\subsection{Spectrum of a compact operator}

Lemma 1. --- Let E be a separated topological vector space of dimension \(\geqslant 2\) over R, and let X be a countable subset of E. The complement E - X is connected.

First, suppose that E has dimension 2. We may assume that \(E = R^{2}\) equipped with the Euclidean norm (EVT, I, p. 14, th. 2). Since X is countable, there exists a real number \(r \in R_{+}^{*}\) such that the circle C with center 0 and radius r does not meet X. Let \(x \in E - X\) ; if \(x \notin C\) , there exists a point \(y \in C\) such that the line \(L_{x}\) joining x and y does not meet X, since X is countable. The set E−X is the union of C, which is connected, and of the connected sets \(L_{x} \cup C\) for \(x \in \mathrm{E} - (\mathrm{X} \cup \mathrm{C})\) ; these sets all contain C, and therefore E-X is connected (TG, I, p. 81, prop. 2).

Consider the general case. By replacing X with X - x for an element \(x \in E - X\) we reduce to the case where \(0 \notin X\) . Since E-X is the union of the sets \(F - (X \cap F)\) as F ranges over the set of 2-dimensional subspaces of E, and since these are connected by the preceding case and contain 0, the set E-X is connected (loc. cit.).

Lemma 2. --- Let S \(\subset \mathbf{C}\) be an infinite, discrete, bounded, and closed set in \(\mathbf{C} - \{0\}\) . Then S is the set of values of a sequence ( \(\lambda_{n}\) ) \(_{n\in N}\) of nonzero complex numbers, pairwise distinct, such that the sequence ( \(|\lambda_{n}|\) ) \(_{n\in N}\) is decreasing and converges to 0.

For every integer \(i \geqslant 1\) , the set \(\mathbf{A}_i\) of complex numbers \(\lambda \in \mathrm{S}\) such that \(|\lambda| \geqslant \frac{1}{i}\) is compact and discrete in \(\mathbf{C}\) , hence finite. Denote its cardinal by \(a_i\) . Since S is infinite, the sequence \((a_i)\) tends to \(+\infty\) . Set \(\mathbf{A}_0 = \varnothing\) and \(a_0 = 0\) . For every \(i \geqslant 1\) , let us choose a bijection \(n \mapsto \lambda_n\) from the interval \([a_{i-1}, a_i[\) of \(\mathbf{N}\) onto \(\mathbf{A}_i - \mathbf{A}_{i-1}\) such that the map \(n \mapsto |\lambda_n|\) is decreasing on \([a_{i-1}, a_i[\) . The sequence \((\lambda_n)_{n \in \mathbf{N}}\) has the required properties.

Let E be a complete normable space of infinite dimension. The algebra \(\mathcal{L}^{\mathrm{c}}(\mathrm{E})\) is a non-unital subalgebra of \(\mathcal{L}(\mathrm{E})\) . Recall that for every compact endomorphism \(u \in \mathcal{L}^{\mathrm{c}}(\mathrm{E})\) the spectrum \(\mathrm{Sp}'_{\mathcal{L}^{\mathrm{c}}(\mathrm{E})}(u)\) is the spectrum of \(u\) relative to the unital subalgebra \(\mathcal{L}^{\mathrm{c}}(\mathrm{E}) \oplus \mathbf{C1}_{\mathrm{E}}\) of \(\mathcal{L}(\mathrm{E})\) (I, p. 4, no. 4).

PROPOSITION 5. — Let E be a complete normable space and let u be a compact endomorphism of E. Every element of \(\mathrm{Sp}_{\mathrm{s}}(u)\) is an eigenvalue of finite spectral multiplicity. Moreover:

\begin{enumerate}
\def\labelenumi{\alph{enumi})}
\item
  If E is finite-dimensional, then \(\operatorname{Sp}(u) = \operatorname{Sp}_{\mathrm{s}}(u)\) ;
\item
  If E is of infinite dimension, then \(\mathrm{Sp}_{\mathrm{s}}(u) = \mathrm{Sp}(u) - \{0\}\) and \(\mathrm{Sp}_{\omega}(u) = \{0\}\) ;
\item
  If \(\mathrm{Sp}_{\mathrm{s}}(u)\) is infinite, it is the set of values of a sequence \((\lambda_{n})_{n\in\mathbf{N}}\) of nonzero complex numbers, pairwise distinct, such that the sequence \((|\lambda_{n}|)_{n\in\mathbf{N}}\) is decreasing and converges to 0;
\item
  If E is of infinite dimension, then \(\mathrm{Sp}(u) = \mathrm{Sp}'_{\mathcal{L}^c (\mathrm{E})}(u)\) .
\end{enumerate}

Every element of \(\mathrm{Sp}_{\mathrm{s}}(u)\) is an eigenvalue of finite spectral multiplicity (prop. 1 of III, p. 83).

The assertion \(a\) ) is elementary (III, p. 85, example 1). Suppose now that E is of infinite dimension.

Let \(\lambda \in \mathrm{Sp}(u) - \{0\}\) . Then \(u - \lambda 1_{\mathrm{E}}\) is a Riesz endomorphism of E (cor. 2 of prop. 2 of III, p. 75), hence \(\lambda\) belongs to \(\mathrm{Sp}_{\mathrm{s}}(u)\) . If E is infinite-dimensional, then \(\mathrm{Sp}_{\omega}(u)\) is nonempty (III, p. 87, th. 1, a)). Necessarily, \(\mathrm{Sp}_{\omega}(u) = \{0\}\) , whence \(b\) ).

The set \(\mathrm{Sp}_{\mathrm{s}}(u)\) is discrete and bounded. Moreover, we have \(\mathrm{Sp}_{\mathrm{s}}(u) = \mathrm{Sp}(u) \cap (\mathbf{C} - \{0\})\) according to \(b\) ), hence \(\mathrm{Sp}_{\mathrm{s}}(u)\) is closed in \(\mathbf{C} - \{0\}\) . The assertion \(c\) ) therefore follows from lemma 2.

The spectrum of \(u\) is countable, and its complement in \(\mathbf{C}\) is therefore connected (lemma 1). By the corollary of prop. 6 of I, p. 28, applied to the unital subalgebra \(\mathcal{L}^{\mathrm{c}}(\mathrm{E})\oplus \mathbf{C}1_{\mathrm{E}}\) of \(\mathcal{L}(\mathrm{E})\) , we conclude that \(\mathrm{Sp}(u) = \mathrm{Sp}'_{\mathcal{L}^{\mathrm{c}}(\mathrm{E})}(u)\) .

PROPOSITION 6. --- Let E be a complete normable space over C and u a compact endomorphism of E.

\begin{enumerate}
\def\labelenumi{\alph{enumi})}
\item
  Let \(f \in \mathcal{O}(\mathrm{Sp}(u))\) such that \(f(0) = 0\) . The endomorphism \(f(u)\) is compact;
\item
  Suppose that E is a complex Hilbert space and that u is normal. Let f be a continuous function on \(\mathrm{Sp}(u)\) such that \(f(0)=0\) . The normal endomorphism \(f(u)\) is compact.
\end{enumerate}

Moreover, the converse assertions are valid if E is infinite-dimensional and the condition \(f(0) = 0\) is not necessary if E is finite-dimensional.

We may assume that E is infinite-dimensional.

Let us prove \(a\) ) and its converse. The endomorphism \(u\) is an element of the Banach algebra \(\mathcal{L}^{\mathrm{c}}(\mathrm{E})\) . Since E is of infinite dimension, we have the equality \(\mathrm{Sp}(u) = \mathrm{Sp}'_{\mathcal{L}^{\mathrm{c}}(\mathrm{E})}(u)\) by prop. 5, \(d\) ). The element \(f(u)\) of the holomorphic functional calculus of the unital Banach algebra derived from \(\mathcal{L}^{\mathrm{c}}(\mathrm{E})\) by adjoining an identity element belongs to \(\mathcal{L}^{\mathrm{c}}(\mathrm{E})\) if and only if \(f(0) = 0\) (I, p. 88). But moreover this element coincides with the element \(f(u)\) of \(\mathcal{L}(\mathrm{E})\) (prop. 7 of I, p. 75), hence \(f(u)\) is compact if and only if \(f(0) = 0\) .

The proof of assertion \(b\) ) and of its converse is entirely identical, considering the continuous functional calculus of the C*-algebra \(\mathcal{L}^{\mathrm{c}}(\mathrm{E})\) (I, p. 110, def. 5).

COROLLARY. --- Let E and F be Hilbert spaces and let u be a continuous linear map from E into F. The linear map u is compact if and only if the endomorphism \(|u|\) of E is compact.

Let \((j, |u|)\) be the polar decomposition of u (def. 4 of I, p. 140). Since \(u = j|u|\) and \(|u| = \sqrt{u^{*}u}\) (prop. 10 of I, p. 139), the equivalence follows from prop. 3 of III, p. 5 and from assertion b) of the preceding proposition.

\subsection{Case of Hilbert spaces}

In this section, E denotes a Hilbert space over C. We denote by \(\pi\) the canonical homomorphism of \(\mathcal{L}(\mathrm{E})\) onto the Calkin algebra \(\mathcal{Calk}(\mathrm{E})\) .

PROPOSITION 7. --- Let \(u \in \mathcal{L}(\mathrm{E})\) .

\begin{enumerate}
\def\labelenumi{\alph{enumi})}
\item
  If u is normal, then u is a Riesz endomorphism if and only if u is a Fredholm endomorphism of index 0.
\item
  If u is Hermitian, then u is a Fredholm endomorphism if and only if u is a Riesz endomorphism.
\end{enumerate}

Every Riesz endomorphism is a Fredholm endomorphism of index 0 (Proposition 5 of III, p. 46). Conversely, suppose that u is a Fredholm endomorphism of index 0, and that u is normal. Its nilspace then coincides with its kernel (EVT, V, p. 43, prop. 8), and is therefore finite-dimensional. It then follows from Lemma 2 of III, p. 45 and from the definition that u is a Riesz endomorphism.

Let us prove \(b\) ). According to \(a\) ), it suffices to check that the index of a Hermitian Fredholm endomorphism \(u\) is zero. But the orthogonal complement of the image of \(u\) (which is closed) is then equal to the kernel of \(u\) (EVT, V, p. 41, prop. 4), whence the assertion.

COROLLARY. --- Let \(u \in \mathcal{L}(\mathrm{E})\) . If u is normal, then \(\mathrm{Sp}_{0}(u)\) is empty, and if u is Hermitian, then \(\mathrm{Sp}_{\mathrm{e}}(u)\) coincides with the spectrum of \(\pi(u)\) relative to the algebra \(\mathcal{Calk}(\mathrm{E})\) .

Both assertions follow from the proposition and the definitions of the sets \(\mathrm{Sp}_{n}(u)\) for \(n \in \overline{Z}\) and of \(\mathrm{Sp}_{\omega}(u)\) , which form a partition of the essential spectrum of u (def. 2 of III, p. 83).

THEOREM 3 (Weyl). --- Let \(u \in \mathcal{L}(\mathrm{E})\) be a normal endomorphism of E. The essential spectrum of \(u\) is the intersection of the sets \(\mathrm{Sp}(u + h)\) for \(h\) ranging over \(\mathcal{L}^{\mathrm{c}}(\mathrm{E})\) .

Let \(h \in \mathcal{L}^c(\mathrm{E})\) . Since \(\mathrm{Sp}_0(u)\) is empty (corollary above), Theorem 2 of III, p. 89 implies that \(\mathrm{Sp_e}(u + h) = \mathrm{Sp_e}(u)\) . The intersection of the sets \(\mathrm{Sp}(u + h)\) therefore contains \(\mathrm{Sp_e}(u)\) .

Let \(\lambda \in \mathrm{Sp}_{\mathrm{s}}(u)\) . Denote by \(\mathrm{E}_{\lambda}\) the eigenspace of \(u\) relative to \(\lambda\) , and \(\mathrm{F}_{\lambda}\) the image of the spectral projector associated with \(u\) and \(\mathbf{C} = \{\lambda\}\) . The space E is the topological direct sum of \(\mathrm{E}_{\lambda}\) and \(\mathrm{F}_{\lambda}\) . Let \(h\) be the finite-rank endomorphism of E that vanishes on \(\mathrm{F}_{\lambda}\) and coincides on \(\mathrm{E}_{\lambda}\) with the identity. The endomorphism \(u + h\) is invertible, hence \(\lambda \notin \mathrm{Sp}(u + h)\) . The theorem follows from this.

Thus, if u is a normal endomorphism of E, and if \(h \in \mathcal{L}^{c}(\mathrm{E})\) , the spectrum of \(u + h\) can differ from the spectrum of u only by isolated points that have finite spectral multiplicity.